\documentclass[14pt,a4paper]{extarticle}

% настройка кодировки, шрифтов и русского языка
\usepackage{fontspec}
\usepackage{polyglossia}

% рабочие ссылки в документе
\usepackage{hyperref}
%\usepackage{cite}

\usepackage[style=numeric, backend=bibtex]{biblatex}
\addbibresource{phys-bib.bib}

\usepackage{csquotes}

%%% Работа с картинками
\usepackage{graphicx} % Для вставки рисунков
\setlength\fboxsep{3pt} % Отступ рамки \fbox{} от рисунка
\setlength\fboxrule{1pt} % Толщина линий рамки \fbox{}
\usepackage{wrapfig} % Обтекание рисунков текстом

% поворот страницы
\usepackage{pdflscape}

% качественные листинги кода
\usepackage{minted}
\usepackage{listings}

% отключение копирования номеров строк из листинга, работает не во всех просмотрщиках (в Adobe Reader работает)
\usepackage{accsupp}
\newcommand\emptyaccsupp[1]{\BeginAccSupp{ActualText={}}#1\EndAccSupp{}}
\let\theHFancyVerbLine\theFancyVerbLine
\def\theFancyVerbLine{\rmfamily\tiny\emptyaccsupp{\arabic{FancyVerbLine}}}

\usepackage{multicol} 
\usepackage{tabularx}

\usepackage{extsizes}
% поля по требованию НИУ ВШЭ / МИЭМ: левое 30 мм, правое 15 мм, верх/низ 20 мм
\usepackage[left=30mm,right=15mm,
top=20mm,bottom=20mm,bindingoffset=0cm]{geometry}

% дополнительные команды для таблиц
\usepackage{booktabs}

% для заголовков
\usepackage{caption} 
\usepackage{titlesec}
\usepackage[dotinlabels]{titletoc}

% разное для математики
\usepackage{amsmath, amsfonts, amssymb, amsthm, mathtools}

% водяной знак на документе, см. main.tex
\usepackage[printwatermark]{xwatermark} 
% настройки polyglossia
\setdefaultlanguage{russian}
\setotherlanguage{english}

% локализация
\addto\captionsrussian{
  \renewcommand{\figurename}{Рисунок}%
  \renewcommand{\partname}{Глава}
  \renewcommand{\contentsname}{\centerline{Содержание}}
  \renewcommand{\listingscaption}{Листинг}
}

% основной шрифт документа
\setmainfont{CMU Serif}

% перечень использованных источников


\usepackage{soul} % Модификаторы начертания
\usepackage{hyperref}
\usepackage[usenames,dvipsnames,svgnames,table,rgb]{xcolor}
\hypersetup{				% Гиперссылки
    unicode=true,           % русские буквы в раздела PDF
    colorlinks=true,       	% false: ссылки в рамках; true: цветные ссылки
    linkcolor=black,          % внутренние ссылки
    citecolor=black,        % на библиографию
    filecolor=magenta,      % на файлы
    urlcolor=blue           % на URL
}
% настройка подсветки кода и окружения для листингов
\usemintedstyle{colorful}
\newenvironment{code}{\captionsetup{type=listing}}{}

% шрифт для листингов с лигатурами
\setmonofont{FiraCode-Regular.otf}[
    Path = templates/,
    Contextuals=Alternate
]

% оформления подписи рисунка
\captionsetup[figure]{labelsep = period}

% подпись таблицы
\DeclareCaptionFormat{hfillstart}{\hfill#1#2#3\par}
\captionsetup[table]{format=hfillstart,labelsep=newline,justification=centering,skip=-10pt,textfont=bf}

% путь к каталогу с рисунками
\graphicspath{{fig/}}

% Внесение titlepage в учёт счётчика страниц
\makeatletter
\renewenvironment{titlepage} {
 \thispagestyle{empty}
}
\makeatother

\counterwithin{figure}{section}
\counterwithin{table}{section}

\titlelabel{\thetitle.\quad}

% для удобного конспектирования математики
\theoremstyle{plain}
\newtheorem{theorem}{Теорема}[section]
\newtheorem{proposition}[theorem]{Утверждение}
\theoremstyle{definition}
\newtheorem{corollary}{Следствие}[theorem]
\newtheorem{problem}{Задача}[section]
\theoremstyle{remark}
\newtheorem*{nonum}{Решение}

% настоящее матожидание
\newcommand{\MExpect}{\mathsf{M}}

% объявили оператор!
\DeclareMathOperator{\sgn}{\mathop{sgn}}

% перенос знаков в формулах (по Львовскому)
\newcommand*{\hm}[1]{#1\nobreak\discretionary{} {\hbox{$\mathsurround=0pt #1$}}{}} 


\begin{document}

% ---- Титульная страница ----
\begin{titlepage}
    \begin{center}
        \large
        Министерство науки и высшего образования\\
        Российской Федерации\\[0.5em]
        Федеральное государственное автономное образовательное\\
        учреждение высшего образования\\[0.5em]
        \textbf{«Национальный исследовательский университет\\
        \mbox{«Высшая школа экономики»}»}\\[1.5em]

        \normalsize
        Московский институт электроники и математики\\
        им.~А.Н.~Тихонова\\[0.3em]
        Департамент прикладной математики\\[2em]

        \large
        Образовательная программа «Прикладная математика»\\
        направление подготовки 01.03.04~«Прикладная математика»\\[2em]

        \textbf{ВЫПУСКНАЯ КВАЛИФИКАЦИОННАЯ РАБОТА}\\[1em]

        на тему\\[0.5em]
        \Large\textbf{«Применение B–LoRA для повышения\\
        качества передачи стиля в генеративных\\
        моделях на базе DiT–FLUX»}\\[2em]

        \large
        \begin{flushright}
            Выполнил студент группы БПМ223:\\
            Бетев Иван Алексеевич\\[1.5em]
            Научный руководитель:\\
            Доцент\\
            Сластников Сергей Александрович\\
        \end{flushright}

        \vfill
        Москва, 2026
    \end{center}
\end{titlepage}

% ---- Задание на ВКР ----
% Задание на выполнение ВКР (вставляется в PDF после титульника)

\newpage
\pagestyle{empty}
{%
    \begin{center}
        \footnotesize
        ПРАВИТЕЛЬСТВО РОССИЙСКОЙ ФЕДЕРАЦИИ\\[0.3em]
        ФЕДЕРАЛЬНОЕ ГОСУДАРСТВЕННОЕ АВТОНОМНОЕ ОБРАЗОВАТЕЛЬНОЕ\\
        УЧРЕЖДЕНИЕ ВЫСШЕГО ОБРАЗОВАНИЯ\\[0.3em]
        \textbf{«НАЦИОНАЛЬНЫЙ ИССЛЕДОВАТЕЛЬСКИЙ УНИВЕРСИТЕТ\\
        «ВЫСШАЯ ШКОЛА ЭКОНОМИКИ»}\\[0.8em]

        \normalsize
        Московский институт электроники и математики\\
        им.~А.Н.~Тихонова\\[2em]

        \Large\textbf{ЗАДАНИЕ}\\[0.3em]
        \large на выполнение выпускной квалификационной работы\\[1.5em]
    \end{center}

    \normalsize
    \noindent
    студенту группы~\textbf{БПМ223}\\[0.4em]
    \textbf{Бетеву Ивану Алексеевичу}\\[1em]

    \noindent\textbf{Тема работы:}\\[0.3em]
    Применение B–LoRA для повышения качества передачи стиля в
    генеративных моделях на базе DiT–FLUX.\\[1em]

    \noindent\textbf{Цель работы:}\\[0.3em]
    Исследовать применимость принципа B-LoRA к архитектуре FLUX.1 для
    управляемого разделения стиля и содержания, экспериментально оценить
    эффективность блочной адаптации относительно существующих подходов и
    установить условия её применимости.\\[1em]

    \noindent\textbf{Формулировка задания:}
    \begin{enumerate}\itemsep0.15em \setlength{\topsep}{0pt}
        \item Систематизировать существующие методы переноса стиля в
              диффузионных моделях (LoRA, IP-Adapter, SplitFlux) и выявить
              их ограничения применительно к FLUX.1.
        \item Сформулировать и экспериментально проверить гипотезу о
              стилевой специализации 19~double-stream блоков FLUX.1,
              реализующих joint attention.
        \item Разработать метод поблочного обучения стилевого
              LoRA-адаптера на double-stream блоках FLUX.1 и схему
              инференса с подключением исключительно стилевого адаптера.
        \item Сформировать экспериментальную базу: обучающие наборы стилей
              Ван~Гога и Моне, оценочные наборы на основе MS-COCO и
              ArtBench-10.
        \item Провести аблационное исследование (диапазон блоков, ранг,
              число шагов, формулировка промпта, коэффициент
              масштабирования) и выбрать оптимальную конфигурацию метода.
        \item Выполнить сравнение предложенного метода с Full-LoRA-FLUX,
              SplitFlux и оригинальным B-LoRA-SDXL по метрикам
              DINO ViT-B/8, CLIP-style, CLIP-content, FID, LPIPS.
        \item Проанализировать режимы отказа метода и сформулировать
              выводы об условиях его применимости.
    \end{enumerate}

    \vspace{1.2em}
    \noindent Проект ВКР должен быть предоставлен студентом в срок до\\[0.2em]
    \mbox{«\,\rule{2em}{0.4pt}\,»\,\rule{8em}{0.4pt}\,20\rule{1.5em}{0.4pt}\,г.}\\[0.4em]
    Научный руководитель ВКР\\[0.2em]
    \mbox{«\,\rule{2em}{0.4pt}\,»\,\rule{8em}{0.4pt}\,20\rule{1.5em}{0.4pt}\,г.}\hfill
    \rule{8em}{0.4pt}\quad С.А.~Сластников\\[0.8em]

    \noindent Первый вариант ВКР предоставлен студентом в срок до\\[0.2em]
    \mbox{«\,\rule{2em}{0.4pt}\,»\,\rule{8em}{0.4pt}\,20\rule{1.5em}{0.4pt}\,г.}\\[0.4em]
    Научный руководитель ВКР\\[0.2em]
    \mbox{«\,\rule{2em}{0.4pt}\,»\,\rule{8em}{0.4pt}\,20\rule{1.5em}{0.4pt}\,г.}\hfill
    \rule{8em}{0.4pt}\quad С.А.~Сластников\\[0.8em]

    \noindent Итоговый вариант ВКР предоставлен студентом в срок до\\[0.2em]
    \mbox{«\,\rule{2em}{0.4pt}\,»\,\rule{8em}{0.4pt}\,20\rule{1.5em}{0.4pt}\,г.}\\[0.4em]
    Научный руководитель ВКР\\[0.2em]
    \mbox{«\,\rule{2em}{0.4pt}\,»\,\rule{8em}{0.4pt}\,20\rule{1.5em}{0.4pt}\,г.}\hfill
    \rule{8em}{0.4pt}\quad С.А.~Сластников\\[1.5em]

    \noindent Задание выдано студенту\\[0.2em]
    \mbox{«\,\rule{2em}{0.4pt}\,»\,\rule{8em}{0.4pt}\,20\rule{1.5em}{0.4pt}\,г.}\hfill
    \rule{8em}{0.4pt}\quad С.А.~Сластников\\
    \hspace*{0pt}\hfill\textit{(научный руководитель ВКР)}\hspace*{1em}\\[0.8em]

    \noindent Задание принято к исполнению студентом\\[0.2em]
    \mbox{«\,\rule{2em}{0.4pt}\,»\,\rule{8em}{0.4pt}\,20\rule{1.5em}{0.4pt}\,г.}\hfill
    \rule{8em}{0.4pt}\quad И.А.~Бетев\\
}
\clearpage
\pagestyle{plain}


% ---- Аннотация (русск.) ----
\section*{Аннотация}
\addcontentsline{toc}{section}{Аннотация}

Данная работа посвящена исследованию применимости принципа блочной
декомпозиции LoRA-адаптеров (метод B-LoRA) к диффузионному трансформеру
FLUX.1 семейства MM-DiT для задачи переноса художественного стиля.
Разработан метод поблочного обучения стилевого адаптера на double-stream
блоках FLUX.1 и предложена методология эмпирической проверки функциональной
специализации блоков на основе техники \textit{prompt injection}. Найдена
оптимальная конфигурация адаптера; проведено сравнение с базовыми подходами
Full-LoRA, SplitFlux и оригинальным B-LoRA-SDXL по метрикам DINO ViT-B/8,
CLIP, FID и LPIPS на наборах стилей Ван~Гога и Моне с использованием
промптов MS-COCO и изображений ArtBench-10.

Полученные результаты подтвердили работоспособность блочной адаптации
в архитектуре MM-DiT и выявили функциональные различия между double-stream
и single-stream блоками FLUX.1, перспективные для дальнейших исследований
эффективного дообучения генеративных моделей.
\newpage

% ---- Abstract (англ.) ----
\section*{Abstract}
\addcontentsline{toc}{section}{Abstract}
\begin{otherlanguage}{english}
This work is devoted to the study of applicability of the block-wise LoRA
decomposition principle (the B-LoRA method) to the FLUX.1 diffusion
transformer of the MM-DiT family for the task of artistic style transfer.
A block-wise training procedure for a style adapter on the double-stream
blocks of FLUX.1 is developed, and a methodology for empirical verification
of the functional specialization of blocks based on the \textit{prompt
injection} technique is proposed. The optimal configuration of the adapter
is established; the proposed method is compared with the baseline approaches
Full-LoRA, SplitFlux and the original B-LoRA-SDXL by DINO ViT-B/8, CLIP,
FID and LPIPS metrics on Van~Gogh and Monet style sets using MS-COCO prompts
and ArtBench-10 images.

The obtained results confirm the feasibility of block-wise adaptation in the
MM-DiT architecture and reveal functional differences between double-stream
and single-stream blocks of FLUX.1, promising for further research on
efficient fine-tuning of generative models.
\end{otherlanguage}
\newpage

% ---- Содержание ----
\tableofcontents
\newpage

% ---- Глава 1: Введение ----
% Введение

\section*{Введение}
\addcontentsline{toc}{section}{Введение}

Актуальность темы исследования определяется стремительным развитием генеративных моделей на основе диффузионного процесса и возросшим практическим интересом к задаче художественного переноса стиля (style transfer). Современные диффузионные модели достигли фотореалистичного качества синтеза изображений, однако управляемое разделение стилистических и семантических характеристик по-прежнему остаётся нетривиальной задачей. Традиционные методы переноса стиля, основанные на свёрточных нейронных сетях~\cite{gatys2016nst}, не позволяют достичь необходимой гибкости при работе с современными трансформерными архитектурами. Появление диффузионных трансформеров (Diffusion Transformer, DiT)~\cite{peebles2023dit} и многомодальных вариантов архитектуры (Multimodal Diffusion Transformer, MM-DiT) открыло новые возможности для исследования внутреннего пространства представлений генеративных моделей.

Модель \textbf{FLUX.1}~\cite{blackforest2024flux}, включающая порядка 12 миллиардов параметров и построенная на базе архитектуры MM-DiT с применением выпрямленных потоков (rectified flow), представляет собой одну из наиболее мощных открытых генеративных систем на сегодняшний день. Её архитектура включает 19 блоков двупоточной обработки (double-stream blocks) и 38 блоков однопоточной обработки (single-stream blocks), что принципиально отличает её от предшествующих моделей на основе U-Net. Несмотря на превосходство FLUX.1 в качестве синтеза, методы эффективной тонкой настройки (fine-tuning) с разделением стиля и содержания для данной архитектуры остаются недостаточно изученными.

Метод \textbf{B-LoRA}~\cite{frenkel2024blora}, предложенный для модели SDXL, продемонстрировал принципиальную возможность разделения стилистических и контентных характеристик посредством избирательного обучения адаптеров низкого ранга (Low-Rank Adaptation, LoRA) в различных блоках U-Net архитектуры. Перенос данного принципа на трансформерную архитектуру FLUX.1 представляет как теоретический, так и практический интерес: MM-DiT обрабатывает текстовые и визуальные токены принципиально иным образом, нежели U-Net, что требует переосмысления гипотезы о специализации блоков.

\textbf{Цель} настоящей работы --- разработать и экспериментально обосновать метод адаптации принципа B-LoRA к архитектуре FLUX.1 для управляемого разделения стиля и содержания при художественном переносе стиля, превосходящий по качеству существующие подходы при сопоставимых вычислительных затратах.

Для достижения поставленной цели необходимо решить следующие \textbf{задачи}:

\begin{enumerate}
    \item Провести систематический анализ существующих методов художественного переноса стиля в диффузионных моделях, включая подходы на основе LoRA, IP-Adapter~\cite{ye2023ipadapter} и конкурирующий метод SplitFlux~\cite{yang2025splitflux}, и выявить их ограничения применительно к архитектуре FLUX.1.
    \item Сформулировать гипотезу о функциональной специализации блоков модели FLUX.1: соответствие блоков двупоточной обработки (блоки 0--5 как начальная гипотеза, уточняемая в ходе аблационного исследования, эксперимент~A) кодированию контентной информации, а блоков однопоточной обработки (блоки 30--37 как начальная гипотеза, также проверяемая экспериментом~A) --- стилистической.
    \item Разработать метод поблочного обучения адаптеров LoRA в рамках архитектуры FLUX.1 с раздельными наборами параметров $\Theta_{\mathrm{content}}$ и $\Theta_{\mathrm{style}}$, предусматривающий совместное обучение обоих адаптеров --- в отличие от метода SplitFlux~\cite{yang2025splitflux}, ограничивающегося исключительно стилевыми блоками, --- а также схему инференса с использованием исключительно стилевого адаптера.
    \item Сформировать экспериментальную базу: подготовить обучающие наборы данных изображений в стиле Ван Гога и Моне, а также оценочные наборы на основе COCO-100~\cite{lin2014coco} и ArtBench-10~\cite{liao2022artbench}.
    \item Провести аблационное исследование (эксперименты A, B, C) для проверки значимости вклада отдельных компонент предложенного метода, включая систематическую проверку гипотезы о выборе диапазонов блоков.
    \item Выполнить сравнительную оценку предложенного метода относительно базовых подходов --- Full-LoRA-FLUX, IP-Adapter-FLUX и B-LoRA-SDXL --- по метрикам CLIP-style, CLIP-content, FID и LPIPS.
    \item Проанализировать полученные результаты и сформулировать выводы об условиях применимости и ограничениях предложенного метода.
\end{enumerate}

\textbf{Объектом исследования} является диффузионная трансформерная модель FLUX.1 как представитель класса архитектур MM-DiT с выпрямленными потоками.

\textbf{Предметом исследования} служат методы поблочной адаптации низкорангового дообучения (LoRA) в рамках архитектуры FLUX.1, обеспечивающие разделение пространств стиля и содержания при художественном переносе стиля.

\textbf{Методы исследования} включают: анализ архитектурных особенностей диффузионных трансформеров; методы параметрически эффективного дообучения (Parameter-Efficient Fine-Tuning, PEFT); количественную оценку качества генерации с использованием мультимодальных метрик на основе модели CLIP~\cite{radford2021clip}; методы статистического сравнения экспериментальных результатов. Реализация выполнена с использованием фреймворков PyTorch~2.1~\cite{paszke2019pytorch}, \texttt{diffusers}~\cite{von2022diffusers}, \texttt{peft}~\cite{peft2022}, системы управления экспериментами ClearML и инструмента \texttt{ai-toolkit}.

\textbf{Научная новизна} работы состоит в следующем: предложен метод совместного поблочного обучения стилевого и контентного адаптеров LoRA ($\Theta_{\mathrm{style}}$ и $\Theta_{\mathrm{content}}$) применительно к архитектуре FLUX.1/MM-DiT; в отличие от метода SplitFlux~\cite{yang2025splitflux}, в настоящей работе оба адаптера обучаются совместно, а роль блоков двупоточной обработки в кодировании контента проверяется систематически; разработана и теоретически обоснована гипотеза о функциональной специализации блоков двупоточной и однопоточной обработки модели FLUX.1, проверяемая посредством систематического аблационного исследования; предложен подход к переносу принципа B-LoRA с архитектуры U-Net на трансформерную архитектуру с сохранением возможности управляемого разделения стиля и содержания.

\textbf{Практическая значимость} работы определяется тем, что разработанный метод позволяет осуществлять высококачественный художественный перенос стиля на основе модели FLUX.1 при использовании малых обучающих наборов данных без модификации весов базовой модели. Предложенный подход применим в задачах цифрового искусства, персонализированной генерации и промышленного дизайна, а также может быть распространён на другие модели класса MM-DiT.

\textbf{Структура работы.} Работа состоит из введения, четырёх глав, заключения и приложений. Глава~2 содержит обзор литературы по методам художественного переноса стиля и параметрически эффективного дообучения диффузионных моделей. В главе~3 описан предложенный метод поблочного обучения адаптеров LoRA в архитектуре FLUX.1. Глава~4 посвящена экспериментальной проверке метода, аблационному исследованию и анализу полученных результатов. В приложении~А приведено описание использованного технологического стека.

\newpage

% ---- Глава 2: Обзор литературы ----
% Глава 2. Обзор существующих методов

\section{Обзор существующих методов}
\label{sec:literature_review}

\subsection{Диффузионные модели генерации изображений}

Диффузионные модели (diffusion models) образуют класс генеративных моделей, основанных на постепенном зашумлении данных и последующем обучении нейронной сети обращать этот процесс. Класс стал одним из ключевых в синтезе изображений, превзойдя генеративно-состязательные сети (generative adversarial networks, GAN)~\cite{dhariwal2021diffbeatgans} по основным показателям качества и разнообразия генерируемых изображений.

Основополагающей для современной эры диффузионных моделей стала работа Хо и соавт.~\cite{ho2020ddpm}, предложивших модель денойзингового диффузионного вероятностного процесса (Denoising Diffusion Probabilistic Model, DDPM). Подход определяет прямой марковский процесс (forward Markov process), в ходе которого к исходному изображению на протяжении $T$ шагов последовательно добавляется гауссовский шум, постепенно превращая его в чистый нормальный шум. Формально, для шага $t$ прямой процесс задаётся как:

\begin{equation}
    q(x_t \mid x_{t-1}) = \mathcal{N}(x_t;\, \sqrt{1-\beta_t}\, x_{t-1},\, \beta_t \mathbf{I})
    \label{eq:ddpm_forward}
\end{equation}

где параметры $\beta_t$ задают расписание шумования (noise schedule). Обратный процесс, реализуемый нейронной сетью, обучается предсказывать добавленный на каждом шаге шум посредством минимизации взвешенного среднеквадратичного отклонения (так называемая $\varepsilon$-параметризация), что эквивалентно денойзинговому согласованию оценок (denoising score matching) при соответствующей перепараметризации~\cite{song2021scoresde}. В качестве архитектуры обратного денойзера использовалась свёрточная сеть в форме U-Net с механизмами внимания. На наборе данных CIFAR-10 модель DDPM достигла показателя FID (Fréchet Inception Distance) 3,17~--- на момент публикации один из наилучших результатов среди likelihood-based генеративных моделей. Работа~\cite{ho2020ddpm} заложила концептуальную и экспериментальную основу всех последующих разработок.

Главное практическое ограничение DDPM~--- высокая вычислительная стоимость генерации: на один образец требуются сотни шагов денойзинга. Сонг и соавт.~\cite{song2021ddim} устранили это ограничение, предложив детерминированный процесс выборки на основе немарковских цепей~--- модель DDIM (Denoising Diffusion Implicit Model). DDIM ускоряет генерацию в 10--50 раз по сравнению с исходным DDPM при сопоставимом качестве. Детерминированность процесса открывает интерполяцию в пространстве скрытых представлений, необходимую для задач редактирования и переноса стиля.

Ромбах и соавт.~\cite{rombach2022ldm} перенесли диффузионный процесс из пространства пикселей в пространство скрытых переменных (latent space) и предложили модель латентной диффузии (Latent Diffusion Model, LDM). Архитектурно LDM состоит из трёх компонентов: вариационного автоэнкодера (Variational Autoencoder, VAE), сжимающего изображение в компактное скрытое пространство; диффузионной модели, работающей непосредственно в этом пространстве; и механизма условного управления через перекрёстное внимание (cross-attention), реализующего текстовое обусловливание на основе языковых эмбеддингов. Это резко снизило вычислительную стоимость обучения и инференса при сохранении качества генерации. На базе архитектуры LDM разработана Stable Diffusion --- открытая модель, ставшая де-факто стандартной платформой для методов управления генерацией, персонализации и переноса стиля. Метод B-LoRA~\cite{frenkel2024blora} изначально разработан на SDXL~--- расширенной версии Stable Diffusion (адаптация к FLUX.1 рассмотрена в разделе~\ref{sec:method}).

Параллельно с развитием диффузионных подходов активно исследовались методы переноса стиля на основе свёрточных нейронных сетей. Работа Гатиса и соавт.~\cite{gatys2016nst} установила формальное разграничение между стилем и содержанием изображения: глубокие признаки VGG-сети, извлечённые из разных слоёв, описывают содержание изображения, тогда как статистика корреляций этих признаков, формализованная через матрицу Грама (Gram matrix), кодирует стилистические характеристики. Метод итеративной оптимизации позволял комбинировать стиль одного изображения с содержанием другого. Несмотря на вычислительную затратность и отсутствие управления через текстовый запрос, этот подход заложил концептуальную основу всех последующих исследований по разделению стиля и содержания, включая B-LoRA.

К 2023 г. диффузионные модели перешли от пиксельного пространства к латентному и от марковских цепей DDPM к детерминированным семплерам DDIM, оформившись как универсальная платформа для синтеза и редактирования изображений с текстовым и визуальным обусловливанием.


\subsection{Диффузионные трансформеры}

По мере роста масштаба диффузионных моделей архитектура U-Net, основанная на свёрточных блоках с остаточными связями и механизмами пространственного внимания, начала проявлять структурные ограничения. Во-первых, свёрточные операции обладают локальным рецептивным полем (receptive field), что затрудняет моделирование дальних пространственных зависимостей. Во-вторых, масштабирование U-Net по числу параметров происходит неоднородно и требует ручной настройки архитектурных гиперпараметров. В-третьих, режим условного обусловливания через перекрёстное внимание в U-Net реализован отдельно от механизма самовнимания, что потенциально ограничивает качество согласования текстового и визуального потоков. Эти ограничения создали предпосылки для перехода к архитектурам на основе трансформеров (transformer-based architectures).

Теоретическую основу всех современных архитектур трансформеров составляет механизм многоголового самовнимания (multi-head self-attention), предложенный в работе Васвани и соавт.~\cite{vaswani2017attention}. Механизм внимания вычисляет взвешенные суммы значений (values) на основе попарного сходства запросов (queries) и ключей (keys):

\begin{equation}
    \text{Attention}(Q, K, V) = \text{softmax}\!\left(\frac{QK^\top}{\sqrt{d_k}}\right) V
    \label{eq:attention}
\end{equation}

Трансформеры устойчиво масштабируются с ростом числа параметров и вычислительных ресурсов (законы масштабирования, scaling laws), что обусловило их выбор как альтернативы свёрточным архитектурам в роли денойзера в диффузионных моделях.

Первое систематическое исследование трансформеров в роли денойзеров диффузионных моделей выполнено Пиблзом и Се~\cite{peebles2023dit} с архитектурой DiT (Diffusion Transformer). Ключевая идея DiT состоит в разбиении скрытого представления на патчи (patches) --- небольшие прямоугольные участки --- и их обработке стандартными блоками трансформера с механизмом самовнимания. Сигнал условного управления, включающий временной шаг диффузии и метку класса, вводится посредством слоёв адаптивной нормализации (adaptive layer normalization, adaLN) с инициализацией нулём (adaLN-Zero), что обеспечивает нулевое воздействие условного сигнала на выход блока в начале обучения и существенно стабилизирует процесс тренировки. На наборе данных ImageNet с разрешением 256×256 пикселей модель DiT достигла значения FID равного 2,27, превзойдя существующие диффузионные модели на основе U-Net. Принципиальный результат: качество модели предсказуемо улучшается с ростом числа параметров и вычислительного бюджета (FLOP), что подтверждает выполнение законов масштабирования для этой архитектуры~\cite{peebles2023dit}.

Эссер и соавт.~\cite{esser2024sd3} разработали мультимодальный DiT (Multimodal DiT, MM-DiT), легший в основу модели Stable Diffusion~3. В отличие от классического DiT, обрабатывающего единый поток токенов, MM-DiT использует архитектуру с двумя параллельными потоками --- отдельными для токенов изображения и токенов текста. Эти потоки взаимодействуют посредством совместного механизма внимания (joint attention), в котором матрицы запросов, ключей и значений для обоих потоков объединяются перед вычислением внимания. Такая архитектура сохраняет за каждым потоком собственные специализированные веса и одновременно обеспечивает двунаправленное взаимодействие визуальных и текстовых представлений. Дополнительно в~\cite{esser2024sd3} предложено улучшенное расписание шумования на основе выпрямленного потока (rectified flow). Совокупность архитектурных решений MM-DiT задала новый уровень качества генерации изображений по текстовому описанию.

Крупнейшая на 2025~г. открытая трансформерная архитектура для генерации изображений~--- FLUX.1, разработанная компанией Black Forest Labs~\cite{blackforest2024flux}. FLUX.1~--- гибридная архитектура, совмещающая блоки MM-DiT с параллельным механизмом внимания (parallel attention) при 12~миллиардах параметров. В качестве текстовых энкодеров используются T5-XXL и CLIP~\cite{radford2021clip}. Архитектура FLUX.1 включает два типа блоков: блоки двойного потока (double stream blocks), реализующие принцип MM-DiT с отдельными потоками для изображений и текста, и блоки единого потока (single stream blocks), обрабатывающие объединённое представление на поздних этапах обработки.

FLUX.1 распространяется в трёх вариантах: \texttt{[pro]}~--- проприетарная версия с наилучшим качеством; \texttt{[dev]}~--- версия с открытыми весами для некоммерческого использования; \texttt{[schnell]}~--- дистиллированная версия под Apache~2.0, генерирующая изображения за 4 шага инференса. По пользовательским предпочтениям FLUX.1 превосходит конкурирующие системы на большинстве оценочных бенчмарков~\cite{blackforest2024flux}.

Переход от U-Net к трансформерным архитектурам меняет организацию информационных потоков: от локальных свёрточных операций~--- к глобальному вниманию, от раздельных потоков текста и изображения~--- к их совместной обработке. Архитектурное разделение на блоки двойного и единого потока в FLUX.1 определяет постановку метода блочно-специфической адаптации в разделе~\ref{sec:method}.


\subsection{Выпрямленные потоки}

Параллельно с развитием архитектур денойзеров улучшался и формализм построения генеративного процесса. Диффузионные модели на марковских цепях имеют ряд недостатков: траектории от шума к данным криволинейны, что требует большого числа шагов для их аппроксимации с приемлемой точностью. Подход на основе потоков (flow-based methods) предлагает альтернативу.

В работе Липмана и соавт.~\cite{lipman2023flow} был предложен метод согласования потока (Flow Matching, FM) --- общий фреймворк для обучения непрерывных нормализующих потоков (Continuous Normalizing Flows, CNF) без использования явного моделирования ОДУ в процессе обучения. Ключевая идея заключается в задании условного вектора скорости, переносящего образцы от шума к данным по оптимальным траекториям. При использовании интерполяции оптимальным транспортом (Optimal Transport, OT) данные траектории оказываются прямолинейными.

Концептуально близкий, но независимо разработанный подход был предложен в работе Лю и соавт.~\cite{liu2022rectifiedflow} под названием выпрямленный поток (Rectified Flow). Выпрямленный поток транспортирует вероятностные меры вдоль прямолинейных траекторий посредством обыкновенных дифференциальных уравнений. Формально, процесс интерполяции задаётся как:

\begin{equation}
    X_t = (1 - t)\, X_0 + t\, X_1, \quad t \in [0, 1]
    \label{eq:rectflow_interp}
\end{equation}

где в обозначениях оригинальной статьи~\cite{liu2022rectifiedflow} $X_0$~--- данные, $X_1$~--- нормальный шум, а нейронная сеть обучается предсказывать постоянную скорость перехода $v = X_1 - X_0$ (в разделе~\ref{sec:method} принята обратная конвенция \texttt{diffusers}~--- содержательный смысл не меняется). Ключевым преимуществом прямолинейных траекторий является возможность использования более редкого шага интегрирования при инференсе: чем прямолинейнее траектория, тем меньше шагов требуется для достаточно точного её следования.

Дополнительно в~\cite{liu2022rectifiedflow} предложена процедура рекурсивного выпрямления (rectification): путём повторного применения операции Reflow на основе ранее сгенерированных пар (данные, шум) можно дополнительно выпрямить траектории, в пределе приближаясь к одношаговой генерации. Формализм выпрямленного потока используется в архитектурах MM-DiT и FLUX.1 в качестве генеративного транспорта. Вариант FLUX.1~\texttt{[schnell]}, способный генерировать изображения за 4~шага инференса, получен из FLUX.1~\texttt{[dev]} не через процедуру рекурсивного Reflow, а через дистилляцию с состязательной целью (adversarial distillation)~\cite{blackforest2024flux}.


\subsection{Методы эффективного дообучения}

Полное дообучение современных диффузионных моделей в миллиарды параметров вычислительно дорого. Для задач персонализации, адаптации стиля и предметно-ориентированной генерации нужны методы, обновляющие лишь небольшую часть весов. Этот класс подходов известен как методы эффективного дообучения с малым числом параметров (Parameter-Efficient Fine-Tuning, PEFT).

Центральный метод этой области~--- низкоранговая адаптация (Low-Rank Adaptation, LoRA), предложенная Ху и соавт.~\cite{hu2022lora}. Ключевая идея: изменения весов предобученной модели при адаптации к новой задаче обладают низким эффективным рангом. Формально, вместо обновления матрицы весов $W \in \mathbb{R}^{d \times k}$ целиком, LoRA параметризует приращение весов как произведение двух матриц малого ранга:

\begin{equation}
    \Delta W = B A, \quad B \in \mathbb{R}^{d \times r},\; A \in \mathbb{R}^{r \times k}, \quad r \ll \min(d, k)
    \label{eq:lora_lit}
\end{equation}

При инференсе приращение $\Delta W$ масштабируется гиперпараметром $\alpha / r$ и складывается с замороженными предобученными весами; при выполнении слияния весов LoRA-адаптера и базовой модели (merge) дополнительная задержка инференса не возникает, тогда как при <<горячем>> подключении адаптера сохраняется небольшое снижение пропускной способности~\cite{hu2022lora}. На практике в диффузионных моделях LoRA обычно применяется к матрицам запросов ($Q$), ключей ($K$), значений ($V$) и выходной проекции в слоях внимания. При ранге $r \in \{4, 8\}$ число обучаемых параметров сокращается в 10~000 раз по сравнению с полным дообучением для моделей масштаба GPT-3 без ощутимой потери качества адаптации~\cite{hu2022lora}.

Из расширений базового подхода выделяется DoRA (Weight-Decomposed Low-Rank Adaptation)~\cite{liu2024dora}, декомпозирующий матрицу весов на компоненты величины (magnitude) и направления (direction) и применяющий LoRA только к компоненте направления, что точнее воспроизводит паттерны полного дообучения. Систематическое сравнение различных вариантов LoRA --- LoHa, LoKr, OFT, BOFT, DoRA --- предпринято авторами библиотеки LyCORIS~\cite{yeh2024lycoris}, предоставляющей комплексный фреймворк оценки применительно к диффузионным моделям.

Применение LoRA к диффузионным моделям для задач персонализации является прямым предшественником метода B-LoRA: если базовый LoRA обрабатывает все блоки модели единообразно, то B-LoRA совершает принципиальный шаг --- ограничивает обучение двумя специфическими блоками, опираясь на гипотезу о структурной специализации блоков применительно к стилю и содержанию.

\begin{figure}[h]
    \centering
    \includegraphics[width=0.98\textwidth]{diffusion_process.png}
    \caption{Концептуальная схема диффузионной генерации и точка
             подключения LoRA-адаптера. Сверху~--- прямой марковский
             процесс зашумления (используется только при обучении);
             снизу~--- обратный процесс, реализуемый нейросетью-денойзером
             ($\varepsilon$-параметризация в DDPM, поле скоростей $v_\theta$
             в формализме rectified flow). Денойзер принимает три условных
             входа: зашумлённый сигнал $x_t$, шаг диффузии $t$ и текстовый
             эмбеддинг. Архитектура денойзера менялась от U-Net (DDPM,
             Stable~Diffusion) к диффузионным трансформерам DiT и MM-DiT
             (FLUX.1, Stable~Diffusion~3). LoRA-адаптация добавляет
             низкоранговые обновления $\Delta W = BA$ к замороженным
             линейным проекциям внимания, не модифицируя сам денойзер.}
    \label{fig:diffusion_process}
\end{figure}


\subsection{Методы персонализации и переноса стиля}

Персонализация генеративных моделей подразумевает адаптацию предобученной модели к воспроизведению специфического визуального концепта по небольшому числу референсных изображений. Перенос стиля~--- частный случай, при котором стилистические характеристики референсного изображения применяются к новому контенту с сохранением его семантики.

Один из первых систематических подходов к персонализации~--- текстовое инвертирование (Textual Inversion)~\cite{gal2023textualinversion}. Метод обучает вектор встраивания для нового псевдослова по 3--5 референсным изображениям при замороженных весах модели. Ключевое достоинство~--- полное сохранение генеративных возможностей базовой модели; ограничение~--- относительно слабая способность воспроизводить детали.

Более мощный подход~--- DreamBooth~\cite{ruiz2023dreambooth}, выполняющий полное дообучение модели на 3--5 изображениях с редким идентификатором-токеном. Для предотвращения семантического дрейфа используется потеря сохранения априорного распределения (prior preservation loss). DreamBooth даёт высокое качество воспроизведения субъекта ценой дорогостоящего полного дообучения.

Иной подход реализован в IP-Adapter~\cite{ye2023ipadapter}: отдельный механизм перекрёстного внимания для признаков изображения-условия обучается при замороженной базовой модели. Адаптер с 22 миллионами параметров проецирует признаки референсного изображения в то же пространство ключей и значений, что и текст, сохраняя полную управляемость через текстовый запрос.

Метод StyleAligned~\cite{hertz2024stylealigned} обеспечивает стилистическую согласованность между несколькими генерируемыми изображениями без обучения: статистики ключей и значений механизма самовнимания разделяются между всеми изображениями батча; применяется адаптивная нормализация по экземпляру (Adaptive Instance Normalization, AdaIN)~\cite{huang2017adain}.

Метод InstantStyle~\cite{wang2024instantstyle} предлагает два решения для улучшения переноса стиля: вычитание признаков содержания из признаков стиля в пространстве IP-Adapter и инъекцию стиля исключительно в блоки перекрёстного внимания, специфичные для стиля. InstantStyle опирается на предположение о блочной специализации, разделяемое методом B-LoRA.

Метод DEADiff~\cite{qi2024deadiff} реализует разделённый подход к кодированию стиля и содержания на архитектуре Q-Former с отдельными токенами-инструкциями. Извлечённые признаки инъецируются в различные слои перекрёстного внимания денойзера. Принципиальная черта~--- использование пар изображений с одинаковым стилем, но разным содержанием, что предотвращает утечку информации между потоками.

Метод B-LoRA~\cite{frenkel2024blora}~--- непосредственный предшественник предлагаемого здесь подхода. Опубликованный на ECCV 2024, B-LoRA строится на эмпирическом наблюдении: в архитектуре SDXL разные блоки кодируют разные аспекты визуальной информации. Идентификация целевых блоков выполнена авторами процедурой инъекции текстовых эмбеддингов (\textit{prompt injection}): подстановка стилевого либо содержательного промпта в каждый блок и анализ отклика генерации позволили установить, что блок \texttt{down\_blocks[2].attentions[1]} в иерархии Down-энкодера (обозначаемый в работе как \textit{W1}) преимущественно отвечает за семантическое содержание изображения, тогда как блок \texttt{up\_blocks[0].attentions[1]} в иерархии Up-декодера (обозначаемый как \textit{W2}) кодирует стилистические характеристики. На основании этого предлагается обучать LoRA-адаптеры исключительно на двух названных блоках, используя в качестве обучающих данных одно-единственное изображение.

Принципиальный результат такого ограниченного обучения~--- неявное разделение стиля и содержания (implicit style/content disentanglement) без явного надзорного сигнала. Обученные адаптеры распадаются на два типа: B-LoRA содержания (Content B-LoRA), ассоциированный с блоком~4, и B-LoRA стиля (Style B-LoRA), ассоциированный с блоком~5. Данные адаптеры могут применяться независимо или совместно, обеспечивая несколько режимов работы: перенос стиля при новом контенте; текстово-управляемая стилизация; генерация изображений, согласованных по стилю. При этом никаких дополнительных шагов обучения для инференса не требуется.

Обучение B-LoRA производится на одном изображении и занимает порядка пяти минут на ускорителе уровня NVIDIA~A100 при числе итераций около~1\,000~\cite{frenkel2024blora}. Потеря обучения является стандартной диффузионной потерей денойзинга применительно к двум целевым блокам. Переобучения не наблюдается благодаря двум факторам: крайне малому числу обучаемых параметров и регуляризирующему эффекту замороженных весов базовой модели, обеспечивающему разнообразие при инференсе.

Экспериментальные результаты~\cite{frenkel2024blora} демонстрируют, что B-LoRA существенно превосходит Textual Inversion~\cite{gal2023textualinversion} и DreamBooth~\cite{ruiz2023dreambooth} по метрикам стилистической верности применительно к задаче переноса стиля в рамках SDXL. Однако B-LoRA имеет принципиальное ограничение: метод разработан и валидирован исключительно на архитектуре SDXL на основе U-Net. Применимость базовой идеи B-LoRA к диффузионным трансформерам~--- центральная проблема настоящей работы.

Магистральное направление развития персонализации и переноса стиля~--- методы, опирающиеся на структурную специализацию блоков модели; B-LoRA наиболее последовательно воплощает эту тенденцию.


\subsection{Адаптация методов переноса стиля к архитектуре FLUX}

С ростом популярности FLUX.1 как генеративной платформы ряд исследований поставил задачу адаптации существующих методов переноса стиля к её специфической архитектуре.

Хуанг и соавт.~\cite{huang2024iclora} демонстрируют способность FLUX.1 к генерации в контексте нескольких изображений (in-context generation): конкатенация нескольких изображений в единое контекстное окно активирует латентные возможности модели, которые затем дообучаются стилезависимым LoRA малого ранга (20--100 образцов). Работа устанавливает совместимость LoRA и FLUX.1, но не исследует специализацию отдельных блоков.

Метод FluxSpace~\cite{dalva2024fluxspace} исследует структуру пространства внимания FLUX.1 с точки зрения редактирования. Авторы показывают, что выходы блоков совместного внимания трансформера образуют семантически значимое пространство представлений, пригодное для управляемого редактирования без обучения на двух уровнях~--- грубой стилистики и тонких атрибутов. Это косвенно свидетельствует о специализации разных блоков FLUX.1 на разные аспекты визуальной информации~--- предположение, лежащее в основе адаптации идей B-LoRA к FLUX.

Наиболее близкая к задаче настоящей работы~--- SplitFlux~\cite{yang2025splitflux}, в которой авторы переносят идею блочной декомпозиции LoRA на архитектуру FLUX.1. Формально следуя гипотезе блочной специализации, SplitFlux расширяет её двумя дополнительными модулями. Авторы устанавливают, что среди блоков единого потока FLUX.1 ранние блоки преимущественно кодируют структурное содержание, тогда как поздние блоки несут ответственность за стилистические характеристики; на основе этого разделения предлагается раздельное обучение LoRA-адаптеров. Помимо базового разделения, SplitFlux вводит два дополнительных механизма: адаптацию с ограничением ранга (Rank-Constrained Adaptation, RCA) для предотвращения деградации идентичности субъекта и визуально-управляемый LoRA (Visual-Gated LoRA Adaptation, VGRA) для фокусировки адаптации на значимых областях.

У SplitFlux есть ряд ограничений. Введение RCA и VGRA усложняет систему по сравнению с концептуальной простотой исходного B-LoRA. Авторы не проводят систематического аблационного анализа роли отдельных блоков. Задействованы только блоки единого потока, тогда как роль блоков двойного потока остаётся неисследованной. Комплексная архитектура затрудняет интерпретацию вклада отдельных компонентов.

Анализ рассмотренных работ выявляет четыре лакуны: отсутствие чистой адаптации B-LoRA к FLUX.1 без дополнительных компонентов; отсутствие систематического анализа блочной специализации FLUX.1; неисследованность роли блоков двойного потока; отсутствие верификации применимости гипотезы B-LoRA при переходе от U-Net к DiT-архитектуре.


\subsection{Выводы по обзору}

Проведённый обзор литературы охватывает ключевые направления исследований по диффузионным моделям, трансформерным архитектурам, выпрямленным потокам, методам PEFT и переносу стиля. Для наглядного сопоставления рассмотренных методов приводится сравнительная таблица~\ref{tab:style_transfer_comparison}.

\begin{landscape}
\begin{table}[p]
\caption{Сравнение методов переноса стиля}
\label{tab:style_transfer_comparison}
\centering
\footnotesize
\begin{tabular}{p{2.8cm} p{2.2cm} p{4.5cm} p{2.5cm} p{2.0cm} p{2.5cm} p{2.0cm}}
\toprule
\textbf{Метод} & \textbf{Базовая модель} & \textbf{Подход} & \textbf{Обучение} & \textbf{Вход стиля} & \textbf{Сохранение содержания} & \textbf{Лицензия} \\
\midrule
Gatys et al.~\cite{gatys2016nst} (2016)
  & VGG
  & Итеративная оптимизация (матрица Грама)
  & Нет (on-image)
  & Изображение
  & Явный content loss
  & Research \\
\addlinespace
Textual Inversion~\cite{gal2023textualinversion} (2023)
  & SD 1.x
  & Оптимизация токен-эмбеддинга
  & Да (per-concept)
  & 3--5 изображений
  & Неявное через промпт
  & CreativeML \\
\addlinespace
DreamBooth~\cite{ruiz2023dreambooth} (2023)
  & SD 1.x/2.x
  & Полное дообучение + prior loss
  & Да (per-subject)
  & 3--5 изображений
  & Prior preservation
  & CreativeML \\
\addlinespace
IP-Adapter~\cite{ye2023ipadapter} (2023)
  & SDXL
  & Разделённый cross-attention адаптер
  & Да (общее)
  & Изображение
  & Сильное
  & Apache 2.0 \\
\addlinespace
StyleAligned~\cite{hertz2024stylealigned} (2024)
  & SDXL
  & Общие attention-статистики + AdaIN
  & Нет
  & Изображение
  & Умеренное
  & Research \\
\addlinespace
InstantStyle~\cite{wang2024instantstyle} (2024)
  & SDXL
  & Блочная инъекция признаков стиля
  & Нет
  & Изображение
  & Сильное
  & Apache 2.0 \\
\addlinespace
DEADiff~\cite{qi2024deadiff} (2024)
  & SDXL
  & Q-Former + нереконструктивное обучение
  & Да (общее)
  & Изображение
  & Сильное
  & Research \\
\addlinespace
\textbf{B-LoRA}~\cite{frenkel2024blora} \textbf{(2024)}
  & \textbf{SDXL}
  & \textbf{Поблочная LoRA, 2 блока, 1 фото}
  & \textbf{Да (${\approx}$5 мин/стиль)}
  & \textbf{1 изображение}
  & \textbf{Сильное}
  & \textbf{Research} \\
\addlinespace
IC-LoRA~\cite{huang2024iclora} (2024)
  & FLUX.1-dev
  & In-context LoRA
  & Да (малый датасет)
  & Набор изображений
  & Задачезависимое
  & Некоммерческая \\
\addlinespace
FluxSpace~\cite{dalva2024fluxspace} (2025)
  & FLUX.1-dev
  & Редактирование в attention-пространстве
  & Нет
  & Изображение
  & Умеренное
  & Research \\
\addlinespace
SplitFlux~\cite{yang2025splitflux} (2025)
  & FLUX.1-dev
  & Поблочная LoRA + RCA + VGRA
  & Да (${\approx}$5 мин/стиль)
  & 1 изображение
  & Сильное (RCA)
  & Research \\
\midrule
\textbf{B-LoRA-FLUX} (предлагаемый)
  & \textbf{FLUX.1-dev}
  & \textbf{Поблочная LoRA на MM-DiT}
  & \textbf{Да (${\approx}$5 мин/стиль)}
  & \textbf{1 изображение}
  & \textbf{Оценивается}
  & \textbf{---} \\
\bottomrule
\end{tabular}
\end{table}
\end{landscape}

Из изложенного складываются четыре магистральные направления развития рассматриваемой области. Во-первых, переход от пиксельных диффузионных моделей к латентным, а затем к архитектурам на основе трансформеров (DiT, MM-DiT, FLUX.1) перестроил организацию информационных потоков в генеративных моделях. Во-вторых, формализм выпрямленных потоков даёт ощутимые практические преимущества над классическими диффузионными процессами. В-третьих, методы PEFT, и LoRA в первую очередь, обеспечивают эффективную адаптацию больших моделей при минимальных вычислительных затратах. В-четвёртых, развитие переноса стиля идёт в сторону методов, опирающихся на структурную специализацию блоков модели.

Анализ выявил четыре конкретные лакуны, обосновывающие постановку задачи настоящего исследования: отсутствие прямого применения идей B-LoRA к FLUX.1; отсутствие систематического изучения специализации блоков FLUX.1; неисследованность роли блоков двойного потока в кодировании стилистической информации; отсутствие верификации применимости гипотезы блочной специализации при переходе от U-Net к DiT-архитектуре.

Перечисленные лакуны определяют цель и методологию работы. Раздел~\ref{sec:chapter3} предлагает методологические решения, адресующие эти пробелы: разработка и обоснование процедуры поиска специализированных блоков в FLUX.1, адаптация обучающей процедуры B-LoRA к трансформерной архитектуре и проектирование схемы инференса для независимого и совместного применения блочно-специфических адаптеров.

\newpage

% ---- Глава 3: Метод и экспериментальная установка ----
\section{Метод и экспериментальная установка}
\label{sec:chapter3}

% Глава 3. Метод и экспериментальная установка
% Подраздел 3.1. Теоретические основы предлагаемого подхода

\subsection{Теоретические основы предлагаемого подхода}
\label{sec:method}

\subsubsection*{Архитектура FLUX.1}

FLUX.1~\cite{blackforest2024flux}~--- семейство диффузионных трансформеров Black~Forest~Labs, построенных на основе многомодального DiT (MM-DiT)~\cite{esser2024sd3,peebles2023dit}. Используется вариант FLUX.1-dev (\texttt{black-forest-labs/FLUX.1-dev}, точная ревизия зафиксирована в конфигурациях; вычисления в bfloat16 на NVIDIA~H100).

Конвейер включает: вариационный автоэнкодер~\cite{rombach2022ldm} со сжатием 8$\times$ в 16-канальное латентное пространство; два текстовых энкодера~--- CLIP~ViT-L/14~\cite{radford2021clip} (77~токенов, pooled projection как условный вектор) и T5-XXL~\cite{raffel2020t5} (до 256~токенов, участвует в joint attention); планировщик \texttt{FlowMatchEulerDiscrete}~\cite{von2022diffusers} при 28~шагах и CFG-guidance~$3{,}5$~\cite{ho2022cfg}. Ключевая особенность~--- разделение бэкбона: первые 19~блоков двухпоточные (double-stream, DS), визуальные и текстовые токены проходят независимыми проекциями с последующим joint attention; следующие 38~блоков однопоточные (single-stream, SS) с конкатенацией модальностей в единый self-attention.

\begin{figure}[h]
    \centering
    \includegraphics[width=0.95\textwidth]{flux_blora_architecture.png}
    \caption{Архитектура FLUX.1-dev и расположение блочных адаптеров B-LoRA. Стилевой адаптер $\Theta_\text{style}$ обучается на ранних double-stream блоках ($[0\text{--}18]$ в основной конфигурации; на схеме акцентированы ds\_00--ds\_01 с наивысшей style-специфичностью по данным Phase~0); содержательный адаптер $\Theta_\text{content}$~--- на поздних double-stream блоках $[9\text{--}18]$. Single-stream блоки в предлагаемом методе не модифицируются.}
    \label{fig:flux_architecture}
\end{figure}

В качестве генеративного транспорта используется метод спрямлённых потоков (rectified flow)~\cite{liu2022rectifiedflow}. В настоящей главе $x_0 \sim \mathcal{N}(0, \mathbf{I})$~--- шум, $x_1$~--- данные, $t \in [0, 1]$~--- время транспорта (нотация противоположна разделу~\ref{sec:literature_review}). Прямолинейный транспорт задаётся ОДУ:
\begin{equation}
    dx = (x_1 - x_0)\,dt,
    \label{eq:rectflow}
\end{equation}
а трансформер обучается предсказывать поле скоростей $(x_1 - x_0)$ минимизацией:
\begin{equation}
    \mathcal{L}(\theta) = \mathbb{E}_{t,\,x_0,\,x_1}\!\left[\bigl\|f_\theta\!\left((1-t)x_0 + t\,x_1,\;t\right) - (x_1 - x_0)\bigr\|^2\right].
    \label{eq:flow_loss}
\end{equation}
Формализм допускает меньшее число шагов интегрирования по сравнению с DDPM и упрощает масштабирование на высокие разрешения.

\subsubsection*{LoRA: адаптация низкого ранга}

LoRA~\cite{hu2022lora} вводит аддитивное обновление низкого ранга для замороженной матрицы весов $W_0 \in \mathbb{R}^{d \times k}$:
\begin{equation}
    W = W_0 + BA, \quad B \in \mathbb{R}^{d \times r},\; A \in \mathbb{R}^{r \times k},\; r \ll \min(d, k).
    \label{eq:lora}
\end{equation}
Полное число обучаемых параметров $r(d+k)$ вместо $dk$: при $r{=}16$ и типичных $d{=}k{=}3072$ для FLUX.1 экономия примерно в 96~раз. Адаптация применяется к проекциям $Q$, $K$, $V$ и \texttt{out\_proj} в каждом блоке. Инициализация $A \sim \mathcal{N}(0, \sigma^2)$, $B = 0$ гарантирует $\Delta W = 0$ в начале обучения. На инференсе вводится множитель силы адаптера $s$:
\begin{equation}
    W = W_0 + s \cdot \frac{\alpha_\text{train}}{r}\, BA,
    \label{eq:lora_scale}
\end{equation}
здесь $\alpha_\text{train}{=}r$ (единичная инициализация по \texttt{peft}); далее $\alpha$ обозначает именно $s$~--- множитель применения адаптера при инференсе.

\subsubsection*{B-LoRA в оригинальной постановке (SDXL)}

Метод B-LoRA~\cite{frenkel2024blora} опирается на эмпирическое наблюдение в SDXL: блок \texttt{up\_blocks[0].attentions[1]} (\textit{W2}, Up-декодер) преимущественно кодирует стиль~--- текстуру, цветовую палитру, художественную манеру,~--- тогда как \texttt{down\_blocks[2].attentions[1]} (\textit{W1}, Down-энкодер) отвечает за структурное содержание. Файнтюн на одном референсе с DreamBooth-промптом~\cite{ruiz2023dreambooth} вида <<\texttt{a [v] in [s] style}>> даёт два независимых адаптера $\Theta_\text{style}$ и $\Theta_\text{content}$, применимых отдельно или совместно. Разделение стиля и содержания достигается имплицитно за счёт структурной специализации блоков, без явных регуляризирующих потерь.

\subsubsection*{Анализ специализации блоков FLUX.1 (Phase~0)}

Поскольку MM-DiT не содержит U-Net-блоков, функциональные роли определяются эмпирически. Предлагается методология \textit{prompt injection}: для каждого блока $b$ текстовый эмбеддинг T5-XXL заменяется эмбеддингом альтернативного промпта, прочие входы остаются неизменными. Степень реакции сгенерированного изображения служит мерой того, насколько блок $b$ «прислушивается» к текстовому условию.

Эксперимент проводится в двух режимах: \textit{style injection} (базовый $P_\text{content}=\texttt{"A photo of a \{object\}"}$, инъецируемый $P_\text{style}=\texttt{"A photo of a \{color\} \{object\}"}$) и \textit{content injection} (направление обратное). $N{=}200$~пар промптов, все 19~DS и все 38~SS-блоки. Мерой реакции служит:
\begin{equation}
    S_b = \frac{1}{N} \sum_{i=1}^{N}
        \cos\!\bigl(\mathrm{CLIP}(I_b^{(i)}),\; \mathrm{CLIP}_\text{text}(P_\text{inject}^{(i)})\bigr),
    \label{eq:phase0_score}
\end{equation}
где $I_b^{(i)}$~--- изображение, сгенерированное с инъекцией в блок~$b$. Высокое $S_b$ указывает на активное участие блока в трансляции текстового условия.

\begin{table}[h]
    \centering
    \small
    \caption{Чувствительность выбранных блоков FLUX.1-dev к инъекции промпта ($N=200$, CLIP ViT-L/14). Специфичность $= S_\text{style} - S_\text{content}$.}
    \label{tab:phase0_results}
    \begin{tabular}{llrrr}
        \toprule
        Блок & Тип & $S_\text{style}$ & $S_\text{content}$ & Специфичность \\
        \midrule
        ds\_00 & DS & 0.260 & 0.243 & \textbf{+0.018} \\
        ds\_01 & DS & 0.229 & 0.225 & +0.004 \\
        ds\_02 & DS & 0.187 & 0.189 & $-$0.003 \\
        ds\_05 & DS & 0.129 & 0.146 & $-$0.017 \\
        ds\_12 & DS & 0.133 & 0.154 & $-$0.021 \\
        ds\_18 & DS & 0.163 & 0.159 & +0.003 \\
        \midrule
        ss\_06 & SS & 0.175 & 0.176 & $-$0.001 \\
        ss\_07 & SS & 0.221 & 0.226 & $-$0.004 \\
        ss\_10 & SS & 0.226 & 0.220 & +0.005 \\
        ss\_23 & SS & 0.242 & 0.240 & +0.002 \\
        ss\_28 & SS & 0.251 & 0.254 & $-$0.002 \\
        ss\_37 & SS & 0.230 & 0.236 & $-$0.006 \\
        \bottomrule
    \end{tabular}
\end{table}

\begin{figure}[h]
    \centering
    \includegraphics[width=\textwidth]{block_analysis_chart.png}
    \caption{Полная карта чувствительности всех 19~double-stream и 38~single-stream блоков FLUX.1-dev по результатам Phase~0 (CLIP-embedding injection, $N=200$~пар промптов). Видны три закономерности: фазовая граница ds\_04$\to$ds\_05, плато single-stream блоков, начиная с ss\_07, и локальные провалы ss\_20/ss\_29.}
    \label{fig:phase0_block_chart}
\end{figure}

\textbf{Граница DS~early\,/\,DS~mid.} Блоки ds\_00 и ds\_01 образуют зону наибольшей style-специфичности; на переходе ds\_04$\to$ds\_05 фиксируется резкий спад $S_\text{style}$ ($\Delta\approx 0.057$). Блоки ds\_05--ds\_12 образуют content-доминирующую долину с экстремумами $-0.021$ (ds\_12) и $-0.017$ (ds\_05).

\textbf{Граница SS~early\,/\,SS~plateau.} Блоки ss\_00--ss\_06 сопоставимы с DS~mid ($S_b\approx 0.17\text{--}0.20$); на переходе ss\_06$\to$ss\_07 происходит скачок $\Delta\approx 0.047$, после чего 31~оставшийся блок образует однородное плато ($S_b\approx 0.22\text{--}0.25$). Style/content-специфичность SS-блоков в плато статистически незначима.

\textbf{Аномальные блоки.} ss\_20 и ss\_29 демонстрируют локальные провалы ($S_b$ ниже соседей на 0.022--0.034) одновременно в обоих режимах, что свидетельствует об архитектурной особенности FLUX.1 и исключает эти блоки из аблационных диапазонов.

\textbf{Выбор блоков.} Кандидаты в стилевые~--- DS~$[0\text{--}18]$ (полный диапазон тестируется в аблации D-A); кандидаты в содержательные (Phase~1b)~--- три перекрывающихся подмножества double-stream: ранний $[0\text{--}8]$ (DC-01), поздний $[9\text{--}18]$ (DC-02) и центральная зона $[5\text{--}12]$ (DC-03). Итоговый выбор $\Theta_\text{content}$ определяется по метрикам обучения.

\subsubsection*{Адаптация B-LoRA к FLUX.1}

Прямой перенос B-LoRA на FLUX.1 невозможен: модель не содержит U-Net и блоков \texttt{up\_blocks}/\texttt{down\_blocks}. Функциональные аналоги стилевых и содержательных блоков ищутся в пространстве MM-DiT экспериментально, а не по архитектурной аналогии.

\textbf{Рабочая гипотеза:} double-stream блоки $[0\text{--}18]$ FLUX.1, реализующие joint attention между текстовым и визуальным потоками, функционально соответствуют стилевому блоку \textit{W2} Up-декодера SDXL~\cite{frenkel2024blora}. Гипотеза согласуется с работами FluxSpace~\cite{dalva2024fluxspace} и SplitFlux~\cite{yang2025splitflux}, показавших функциональную специализацию блоков MM-DiT.

По результатам Phase~1 и Phase~1b в контрольной точке GATE~1 принимается один из трёх сценариев:
\begin{itemize}\itemsep0.2em
    \item \textbf{Сценарий~A}~--- DS-блоки дают перенос стиля \textit{и} выделен $\Theta_\text{content}$: полное разделение стиля и содержания (дополнительно выполняется Phase~6 style-content mixing);
    \item \textbf{Сценарий~B}~--- DS-блоки дают перенос стиля, но $\Theta_\text{content}$ не выделяется: применяется только $\Theta_\text{style}$, Phase~6 пропускается;
    \item \textbf{Сценарий~C}~--- DS-блоки не обеспечивают переноса: проводится дополнительный анализ альтернативных подмножеств, конфигурация переопределяется.
\end{itemize}

\begin{figure}[h]
    \centering
    \includegraphics[width=0.98\textwidth]{gate1_decision.png}
    \caption{Структура решения в контрольной точке GATE~1.
             Phase~1 (серия~D на DS~$[0\text{--}18]$) и Phase~1b (серия~DCC,
             поиск $\Theta_\text{content}$) образуют двумерный критерий:
             значимость стилевого сдвига и наличие выделяемого содержательного
             адаптера. По итогам работы (раздел~\ref{sec:results_phase1b})
             принят Сценарий~A: $\Theta_\text{style}{=}$DS~$[0\text{--}18]$,
             $\Theta_\text{content}{=}$DS~$[9\text{--}18]$.}
    \label{fig:gate1_decision}
\end{figure}

\textbf{Реализация метода.} (1)~Все веса FLUX.1 замораживаются. (2)~К double-stream блокам $[0\text{--}18]$ добавляется LoRA-адаптер ранга~$r$~--- $\Theta_\text{style}$; выделение $\Theta_\text{content}$ производится при Сценарии~A. (3)~Адаптер обучается на одном референсе за $T$~шагов (из $\{500,\,1\,000,\,2\,000,\,4\,000\}$, оптимум по~\ref{sec:ablations}) с целевой функцией flow-matching~\eqref{eq:flow_loss}, скоростью обучения $5{\times}10^{-5}$, оптимизатором AdamW~\cite{loshchilov2019adamw} ($\beta_1{=}0{,}9$, $\beta_2{=}0{,}999$, weight decay $10^{-2}$, без warmup, без EMA), в bfloat16 средствами Hugging~Face \texttt{accelerate}~\cite{accelerate2022} и фреймворка \texttt{ai-toolkit}. (4)~При инференсе подключается $\Theta_\text{style}$: модель воспроизводит стиль референса, а содержание задаётся текстовым промптом.

\subsubsection*{Поиск адаптера содержания $\Theta_\text{content}$ (Phase~1b)}

В оригинальном B-LoRA~\cite{frenkel2024blora} $\Theta_\text{style}$ и $\Theta_\text{content}$ приписаны заранее известным блокам SDXL; в FLUX.1 такого априорного соответствия нет. По данным Phase~0 наибольшая content-специфичность сосредоточена в диапазоне ds\_05--ds\_12. В Phase~1b обучаются и сравниваются три подмножества: ранний $[0\text{--}8]$ (DC-01), поздний $[9\text{--}18]$ (DC-02), центральный $[5\text{--}12]$ (DC-03). На каждом диапазоне обучается отдельный LoRA-адаптер ($r{=}16$, прочие гиперпараметры как в Phase~1); при инференсе применяется только этот адаптер без $\Theta_\text{style}$.

Адаптер признаётся содержательным, если относительно базовой модели достигается (i)~рост CLIP-content или (ii)~статистически значимое отклонение DINO-style с ограничением подъёма фактором~$<\!5$. При выполнении условия хотя бы для одного диапазона фиксируется \textbf{Сценарий~A}, иначе~--- \textbf{Сценарий~B}; обработка \textbf{Сценария~C} описана выше.


% Раздел: Используемые данные (глава «Метод и экспериментальная установка»)

\subsection{Используемые данные}
\label{sec:datasets}

B-LoRA относится к однофотографическим подходам: для файнтюна достаточно одного изображения-референса. Экспериментальная оценка дополнительно требует референсного распределения для FID и набора контентных промптов для семантической оценки. Сводка наборов приведена в таблице~\ref{tab:datasets}.

\textbf{Эталонные стили (DS1).} Отобрано 8~изображений двух художественных школ~--- импрессионизма Моне и постимпрессионизма Ван~Гога,~--- получены с WikiArt и выбраны по критериям стилистической выраженности и цветовой характерности. Для Ван~Гога: «Звёздная ночь» (1889), «Красные виноградники» (1888), «Подсолнухи» (1888), «Автопортрет с перевязанным ухом» (1889); для Моне: «Стог сена в~Живерни» (1886), «Водяные лилии~--- вечерний эффект» (1899), «Мост Ватерлоо» (1903), «Мадам Моне с ребёнком» (ок.~1875). При каждом запуске обучения используется одно изображение; набор из четырёх позволяет оценить вариативность результатов внутри стиля. Выбор обусловлен стилистической контрастностью: импрессионизм Моне отличается размытыми контурами и пастельной палитрой, манера Ван~Гога~--- вихреобразной фактурой и насыщенным цветом.

\begin{figure}[h]
    \centering
    \includegraphics[width=0.95\textwidth]{style_refs_grid.png}
    \caption{Стилевые референсы. Верхний ряд~--- Ван~Гог: «Красные виноградники в~Арле» (1888), «Автопортрет с~перевязанным ухом» (1889), «Подсолнухи» (1888), «Звёздная ночь» (1889). Нижний ряд~--- Моне: «Стог сена в~Живерни» (1886), «Мадам Моне с~ребёнком» (ок.~1875), «Водяные лилии~--- вечерний эффект» (1899), «Мост Ватерлоо» (1903). Источник: WikiArt; работы в общественном достоянии.}
    \label{fig:style_refs}
\end{figure}

\textbf{ArtBench-10 (DS3).} Датасет~\cite{liao2022artbench} (около 60\,000~изображений по 10~стилям) используется для оценки FID в \textit{reduced protocol}: подвыборка по 500~изображений на жанр. Из-за ограниченного объёма метрика применяется для сравнения конфигураций между собой, а не для сопоставления с опубликованными значениями. Источник: HuggingFace Hub (\texttt{Doub7e/ArtBench-10}), CC-BY.

\textbf{MS-COCO val2017 (DS2).} Подвыборка из 100~промптов валидационного раздела~\cite{lin2014coco} обеспечивает нейтральный контентный фон для измерения CLIP-сходств. Источник: \url{cocodataset.org}, CC-BY.

\textbf{DreamBooth (DS5).} Репозиторий \texttt{google/dreambooth}~\cite{ruiz2023dreambooth}, Apache~2.0. Используются 8~предметов: \texttt{dog} ($\approx$30~изображений~--- основной объект Phase~1b) и 7~дополнительных (\texttt{cat}, \texttt{backpack}, \texttt{bowl}, \texttt{can}, \texttt{clock}, \texttt{vase}, \texttt{bear}, по~$\approx$25~изображений). Совокупный пул~--- $\approx$200~контентных изображений для количественного сравнения методов (Phase~5.2).

\textbf{Парный набор DS8.} Воспроизводит протокол раздела~5.1 оригинальной статьи~\cite{frenkel2024blora} в виде \textit{reduced protocol}: $N_c{=}8$ объектов из DS5, $N_s{=}8$ стилей из DS1; из декартова произведения $8{\times}8{=}64$ семплируется 50~пар при \texttt{seed=42}. Манифест \texttt{experiments/data/b\_lora\_eval\_pairs.json} версионирован через DVC. Основная метрика~--- среднее косинусное сходство DINO~ViT-B/8 между сгенерированным изображением и стилевым референсом. Поскольку пул отличается от оригинального, абсолютные значения не сопоставимы напрямую с опубликованной Table~1; относительные различия между методами остаются корректными.

\begin{table}[h]
    \centering
    \caption{Сводка использованных данных}
    \label{tab:datasets}
    \begin{tabular}{p{2.8cm}p{4.2cm}p{4.0cm}p{3.4cm}}
        \hline
        \textbf{Датасет} & \textbf{Назначение} & \textbf{Размер} & \textbf{Лицензия} \\
        \hline
        WikiArt (DS1) & Обучение B-LoRA; DS8 стили & 8~изображений (2~стиля~$\times$~4) & общественное достояние \\
        DreamBooth~\cite{ruiz2023dreambooth} (DS5) & DS8 контент; Phase~1b & 8~предметов ($\approx$200~изображений) & Apache 2.0 \\
        ArtBench-10~\cite{liao2022artbench} (DS3) & Референс FID (approximate) & 500~изображений на жанр (reduced) & CC-BY \\
        MS-COCO val2017~\cite{lin2014coco} (DS2) & Инференс; CLIP-content & 100~промптов (подвыборка) & CC-BY \\
        DS8: парный eval & Сравнение методов (Phase~5.2) & 8~$\times$~8~$\to$~50~пар (reduced) & DS1~+~DS5 \\
        \hline
    \end{tabular}
\end{table}

Все наборы версионированы средствами DVC: \texttt{.dvc}-файлы фиксируют хэши содержимого в git-репозитории, а сами данные хранятся в S3-совместимом удалённом хранилище. Команда \texttt{dvc pull} восстанавливает идентичный набор после клонирования.


% Раздел: Подготовка данных (глава «Метод и экспериментальная установка»)

\subsection{Подготовка данных}
\label{sec:preprocessing}

\textbf{Стилевые изображения.} Восемь эталонных работ организованы в файловую иерархию \texttt{data/styles/\{van\_gogh,monet\}/img\{1..4\}/}, совместимую с параметрами \texttt{folder\_path} и \texttt{caption\_ext} фреймворка \texttt{ai-toolkit}: каждое изображение сопровождается одноимённым \texttt{.txt}-файлом с подписью. Использование четырёх изображений на художника соответствует постановке оригинального B-LoRA~\cite{frenkel2024blora} (для файнтюна достаточно одного референса) и позволяет оценить вариативность результатов внутри стиля при сопоставимом вычислительном бюджете.

Подписи составлены вручную и описывают \textit{исключительно} семантическое содержание~--- предметы, пространственные отношения, освещение~--- без упоминаний имени художника или принадлежности к стилю. Это требование принципиально для B-LoRA: стилевая информация должна кодироваться в весах $\Theta_\text{style}$, а не передаваться через текстовое условие. Например, подпись к «Звёздной ночи»~--- \textit{<<a night sky with swirling clouds and stars over a village with a church steeple>>}; к «Красным виноградникам»~--- \textit{<<a vineyard at sunset with red fields and workers in a rural landscape>>}. В конфигурации обучения задан \texttt{caption\_dropout\_rate: 0.05}: на 5\,\% шагов обучение проводится без текстового условия~\cite{ho2022cfg}, что снижает зависимость адаптера от конкретных формулировок и повышает обобщающую способность.

\textbf{Промпты для оценки.} Сформирован набор из 100~текстовых промптов, извлечённых из \texttt{captions\_val2017.json}: \texttt{seed=42}, длина 5--20~слов, без дубликатов, сохранение в \texttt{data/coco\_prompts.txt}. Размер выбран как компромисс между вычислительной стоимостью генерации (100~прогонов FLUX.1 на эксперимент при 28~шагах инференса) и статистической надёжностью агрегированных метрик; требования~--- семантическое разнообразие сцен и нейтральность по отношению к художественному стилю~\cite{lin2014coco}.

\textbf{Ограничение FID.} Надёжная оценка FID требует тысяч изображений~\cite{heusel2017fid,parmar2022cleanfid}; в работе используется 100~сгенерированных против 500~реальных из ArtBench-10. Вычисляемое значение носит ориентировочный характер и используется для сравнения методов в рамках одного протокола, а не как абсолютная оценка качества.


% Раздел: Протокол оценки (глава «Метод и экспериментальная установка»)

\subsection{Протокол оценки}
\label{sec:eval_protocol}

\subsubsection*{Метрики качества}

Для количественной оценки результатов переноса стиля применяется набор из пяти
метрик, каждая из которых измеряет отдельный аспект качества генерации.

\textbf{DINO-style}~--- косинусное сходство между \texttt{[CLS]}-токенами модели
DINO ViT-B/8~\cite{caron2021dino} для сгенерированного изображения и стилевого
референса. Эта метрика является \textit{основным} критерием оценки качества переноса
стиля: авторы B-LoRA~\cite{frenkel2024blora} используют именно DINO-сходство в качестве
ключевой метрики (Table~1 оригинальной статьи), поскольку DINO-признаки чувствительны
к текстуре и художественной манере независимо от семантики сцены. Вычисляется по модели
\texttt{facebook/dino-vitb8} (Hugging Face \texttt{transformers}~\cite{wolf2020transformers})
как среднее по батчу косинусное сходство нормализованных \texttt{[CLS]}-эмбеддингов.

\textbf{CLIP-style}~--- косинусное сходство между эмбеддингами CLIP~\cite{radford2021clip}
сгенерированного изображения и стилевого референса. Метрика оценивает степень
соответствия визуального стиля эталону: чем выше значение, тем точнее воспроизведена
художественная манера. Вычисляется по модели CLIP ViT-L/14 на нормализованных векторах
с использованием L2-нормировки перед скалярным произведением; абсолютные значения CLIP-style
в задаче переноса стиля ограничены конструкцией CLIP-пространства, однако при значениях
$\gtrsim 0{,}7$ интерпретация метрики становится осторожной (см. ограничения в гл.~<<Результаты и анализ>>).

\textbf{CLIP-content}~--- косинусное сходство между эмбеддингом CLIP сгенерированного
изображения и текстовым эмбеддингом входного промпта. Метрика оценивает, насколько
хорошо сохраняется семантическое содержание при переносе стиля: высокое значение
означает, что стилизованное изображение соответствует заданному промпту.

\textbf{FID} (\textit{Fréchet Inception Distance})~\cite{heusel2017fid}~--- расстояние
Фреше между распределениями активаций сети Inception-v3 на реальных и сгенерированных
изображениях. Вычисляется с использованием библиотеки clean-fid~\cite{parmar2022cleanfid}
на выборке из датасета ArtBench-10 (500~реальных изображений соответствующего жанра против
100~сгенерированных). Низкое значение FID соответствует высокому фотореализму и близости
к целевому стилевому распределению.

\textbf{LPIPS} (\textit{Learned Perceptual Image Patch Similarity})~\cite{zhang2018lpips}~---
перцептивное сходство, вычисляемое на основе активаций сети VGG. Используется как
вспомогательная метрика для оценки текстурного соответствия сгенерированного изображения
и стилевого референса. Низкое значение LPIPS соответствует высокому перцептивному
сходству (метрика расстояния).

\begin{figure}[h]
    \centering
    \includegraphics[width=\textwidth]{eval_pipeline.png}
    \caption{Схема пайплайна оценки. Верхний ряд~--- линейная последовательность
             обучения и инференса: одно стилевое изображение DS1 используется
             для обучения LoRA-адаптера $\Theta_\text{style}$; на этапе
             инференса 100 промптов MS-COCO val2017 (DS2) пропускаются через
             FLUX.1-dev с подключённым адаптером и дают 100~стилизованных
             изображений. Нижний ряд~--- пять метрик, вычисляемых на полученной
             выборке. Стрелки сбоку показывают вспомогательные входы: стилевой
             референс участвует в~DINO-style и CLIP-style; ArtBench-10 (DS3)
             служит реальной выборкой для FID. Все генерации детерминированы
             ($\mathtt{seed}=42$); агрегация выполняется усреднением по~100~промптам и
             4~опорным изображениям внутри стиля.}
    \label{fig:eval_pipeline}
\end{figure}

\subsubsection*{Схема оценки}

Полный протокол включает на один метод
$N_\text{gen}^{(1)} = N_\text{style} \times N_\text{ref}\times N_\text{prompts}
= 2 \times 4 \times 100 = 800$~генераций (два стиля Van~Gogh и Monet, по
4~стилевых референса в каждом, 100~промптов MS-COCO val2017); для основного
сравнения трёх методов (Phase~4) суммарно $N_\text{gen}=N_\text{methods}\cdot N_\text{gen}^{(1)}=2\,400$~генераций. Дополнительно
$1\,000$~генераций для оценки FID на подмножестве ArtBench-10
($100~\text{промптов} \times 10~\text{стилей ArtBench}$, seed фиксирован = 42).

Для каждого метода:
\begin{enumerate}
    \item Обучить LoRA-адаптер $\Theta_\text{style}$ на одном стилевом изображении
          (по одному прогону на каждое из 4~изображений стиля).
    \item Сгенерировать изображения по 100~промптам из MS-COCO val2017
          с применением $\Theta_\text{style}$.
    \item Вычислить DINO-style (косинусное сходство \texttt{[CLS]}-токенов DINO ViT-B/8
          между сгенерированными изображениями и стилевым референсом) — основная метрика.
    \item Вычислить CLIP-style (сравнение с референсным изображением)
          и CLIP-content (сравнение с текстовым промптом).
    \item Вычислить FID между 1000~сгенерированными изображениями
          и 500~изображениями ArtBench-10 соответствующего жанра
          (\textit{approximate FID}, reduced protocol).
\end{enumerate}

Итоговые значения метрик усредняются по 4~стилевым изображениям внутри каждого стиля
(Van~Gogh, Monet). Все генерации производятся при фиксированном значении seed~=~42,
шагах сэмплирования~=~28 и guidance scale~=~3{,}5.

\subsubsection*{Сравниваемые методы}

В основном эксперименте (Phase~4) сравниваются четыре метода:

\begin{itemize}
    \item \textbf{B-LoRA-FLUX} (предлагаемый) — LoRA-адаптер на double-stream
          блоках FLUX.1, конфигурация которых (диапазон блоков, ранг~$r$, число шагов)
          определяется по результатам аблационного исследования (раздел~\ref{sec:ablations});
          при инференсе применяется только $\Theta_\text{style}$.
    \item \textbf{Full-LoRA-FLUX} — LoRA-адаптер на всех блоках FLUX.1 (без
          блочного разделения); служит для проверки гипотезы о необходимости
          селективного выбора блоков.
    \item \textbf{SplitFlux}~\cite{yang2025splitflux} --- метод блочного разделения
          стиля и содержания в FLUX.1, основанный на single-stream блоках; выступает
          в роли конкурирующего метода внутри архитектуры FLUX и позволяет оценить
          относительный вклад выбора типа блоков (double-stream vs.\ single-stream).
    \item \textbf{IP-Adapter-FLUX}~\cite{ye2023ipadapter}~--- метод на основе
          адаптера изображения, не требующий файнтюнинга. На момент проведения
          экспериментов в открытом доступе отсутствовала стабильная публичная
          реализация IP-Adapter, совместимая с актуальной версией FLUX.1-dev:
          в ходе работы было опробовано несколько кандидатных репозиториев,
          однако ни один из них не обеспечил воспроизводимого пайплайна с
          использованной версией модели и весами. По этой причине IP-Adapter
          не включён в итоговое сравнение, и данный пробел зафиксирован как
          ограничение работы. Повторение сравнения после появления стабильной
          реализации указано в перспективах дальнейших исследований.
\end{itemize}

Дополнительно в Phase~5 проводится кросс-архитектурное сравнение с
\textbf{B-LoRA-SDXL}~\cite{frenkel2024blora} --- оригинальным методом на базовой
модели SDXL --- для оценки прироста качества при переходе к архитектуре FLUX.1
(подробнее~--- параграф «Phase~5» ниже).

\subsubsection*{Аблационные исследования}
\label{sec:ablations}

Аблационное исследование проводится в несколько последовательных фаз на стиле Van~Gogh
(одно фиксированное изображение, seed~=~42). Фазы выполняются в строго определённом
порядке; результаты каждой фазы определяют конфигурацию следующей. Переход между
фазами регулируется контрольной точкой GATE~1, фиксирующей итоговое решение о выборе
целевых блоков и сценарии дальнейшей работы.

\paragraph{Анализ блоков FLUX.1 (Phase~0).}
Перед обучением каких-либо адаптеров проводится анализ специализации блоков методом
prompt injection (методология описана в разделе~\ref{sec:method}, уравнение~\eqref{eq:phase0_score}).
Для каждого из 19~double-stream блоков и подмножества single-stream блоков
$\{0, 10, 20, 30, 37\}$ вычисляется оценка $S_b$ в режимах style injection
и content injection на 200~парах промптов. Результаты Phase~0 используются для:
\begin{enumerate}
    \item определения приоритетных диапазонов блоков для аблации~D-A;
    \item выявления кандидатов на роль содержательных блоков для Phase~1b.
\end{enumerate}

\paragraph{Диагностическая серия D (Phase~1).}
Базовая проверка рабочей гипотезы. Включает четыре конфигурации:

\begin{itemize}
    \item \textbf{e00\_no\_lora\_baseline} — генерация без LoRA-адаптера; служит
          нижней границей по DINO-style и CLIP-style.
    \item \textbf{d01} — double-stream блоки~[0–18], $r{=}16$, 1\,000~шагов,
          $\text{lr}{=}5{\times}10^{-5}$, промпт \texttt{"a sks"}
          (минималистичный аналог \texttt{"A [v]"} из оригинальной статьи~\cite{frenkel2024blora}).
    \item \textbf{d02} — те же блоки и ранг, 2\,000~шагов (проверка влияния
          числа итераций дообучения).
    \item \textbf{d03} — double-stream блоки~[0–18], $r{=}32$, 1\,000~шагов
          (проверка влияния ранга).
\end{itemize}

\paragraph{Поиск адаптера содержания (Phase~1b).}
По результатам Phase~0 на трёх перекрывающихся диапазонах double-stream блоков~---
ранних $[0\text{--}8]$ (конфигурация DC-01), поздних $[9\text{--}18]$ (DC-02) и
центральной зоне $[5\text{--}12]$, гипотетически предпочтительной по Phase~0 (DC-03),~---
обучается отдельный LoRA-адаптер ($r{=}16$, те же гиперпараметры, что в~d01).
При инференсе применяется только данный адаптер; измеряются CLIP-content и DINO-style.
Адаптер признаётся содержательным ($\Theta_\text{content}$), если CLIP-content
возрастает относительно базовой линии, а DINO-style остаётся ограниченным
(см.\ критерии в разделе~\ref{sec:method}).

\paragraph{Контрольная точка GATE~1.}
По совокупным результатам Phase~1 и Phase~1b принимается одно из трёх решений:
\begin{itemize}
    \item \textbf{Сценарий~A} --- double-stream блоки обеспечивают перенос стиля
          (DINO-style в~d01 существенно выше базовой линии) \textit{и} выделен
          $\Theta_\text{content}$: выполняются полные аблации~D-A–D-P, Phase~3,
          Phase~3.2, Phase~4, Phase~4b, Phase~5, Phase~6.
    \item \textbf{Сценарий~B} --- double-stream блоки дают перенос стиля, но
          $\Theta_\text{content}$ не выделяется: выполняются аблации~D-A–D-P,
          Phase~3, Phase~4, Phase~4b, Phase~5; Phase~3.2 и Phase~6 пропускаются.
    \item \textbf{Сценарий~C} --- double-stream блоки не дают видимого переноса
          стиля: проводится дополнительный анализ альтернативных подмножеств
          блоков; конфигурация переопределяется до начала аблаций.
\end{itemize}

\paragraph{Аблация D-A: диапазон double-stream блоков (Phase~2.1).}
При фиксированном ранге и числе шагов (из лучшей конфигурации серии~D)
проверяются четыре диапазона double-stream блоков: [0–6], [0–12], [0–18], [6–18].
Цель: установить, какой диапазон блоков обеспечивает наилучший перенос стиля
по метрике DINO-style при сохранении семантической связи с промптом (CLIP-content).

\paragraph{Аблация D-B: ранг LoRA (Phase~2.2).}
При фиксированном диапазоне блоков (лучший по D-A) проверяются значения
$r \in \{4, 16, 32, 64\}$. Значение $r{=}64$ соответствует конфигурации
оригинальной статьи B-LoRA~\cite{frenkel2024blora} и включается для прямой
сопоставимости результатов.

\paragraph{Аблация D-C: число шагов (Phase~2.3).}
При фиксированных блоках и ранге (лучшие по D-A и D-B) проверяются значения
$\text{steps} \in \{500, 1\,000, 2\,000, 4\,000\}$. Цель: определить компромисс
между качеством переноса стиля и временем дообучения.

\paragraph{Аблация D-P: текстовый промпт обучения (Phase~2.4).}
При финальной конфигурации блоков, ранга и числа шагов (лучшие по D-A–D-C)
проверяются четыре варианта обучающего промпта:
\texttt{"a sks"}, \texttt{"a sks painting"},
\texttt{"a [v]"}, \texttt{"a [v] painting in [s] style"}.
Цель: оценить влияние текстового кондиционирования при обучении на степень
разделения стиля и содержания. Метрики: DINO-style, CLIP-content, CLIP-style
на 100~генерациях (промпты из MS-COCO val2017).

\paragraph{Аблация G: коэффициент применения адаптера (Phase~3).}
При финальной конфигурации блоков, ранга, числа шагов и промпта (лучшие по D-A–D-P)
проверяются значения lora\_scale $\in \{0{,}3;\; 0{,}5;\; 0{,}7;\; 1{,}0;\; 1{,}5;\; 2{,}0\}$.
Согласно оригинальной статье~\cite{frenkel2024blora}, оптимальное значение
для переноса стиля лежит в диапазоне~$\alpha \approx 0{,}4{-}0{,}5$.

\paragraph{Раздельные коэффициенты $\alpha_\text{style}$ и $\alpha_\text{content}$ (Phase~3.2, только Сценарий~A).}
При наличии двух независимых адаптеров ($\Theta_\text{style}$ и $\Theta_\text{content}$)
исследуется влияние их совместного масштабирования. Тестируются следующие комбинации:
$\alpha_\text{style} \in \{0{,}4;\; 0{,}5;\; 0{,}7;\; 1{,}0\}$ при фиксированном
$\alpha_\text{content} = 1{,}0$ и, обратно, $\alpha_\text{content} \in \{0{,}4;\; 0{,}5;\; 0{,}7;\; 1{,}0\}$
при фиксированном $\alpha_\text{style} = 1{,}0$. Метрики: DINO-style, CLIP-content,
CLIP-style на 100~генерациях. Цель: найти баланс между силой стилизации
и сохранением семантики содержания.

По совокупным результатам аблаций D-A–D-C, D-P и~G выбирается итоговая конфигурация
с наилучшим балансом DINO-style и CLIP-content; она применяется в основном
эксперименте Phase~4 и финальной проверке на стиле Monet.

\paragraph{Анализ ограничений (Phase~4b).}
После основного сравнения методов (Phase~4) проводится систематический анализ
типичных артефактов генерации. Рассматриваются три класса ограничений:
\begin{itemize}
    \item \textbf{Color leakage} --- проникновение цветовой палитры стилевого
          референса в области изображения, где оно нежелательно. Количественная
          оценка: средняя цветовая разница $\Delta E$ (CIELAB) между генерацией
          и нейтральным baseline по 5~промптам с нейтральным фоном, 5~генераций
          на промпт.
    \item \textbf{Background leakage} --- сохранение структурных элементов фона
          стилевого референса в генерациях. Оценка: DINO-object similarity
          (косинусное сходство DINO ViT-B/8 на патчах переднего плана относительно
          фона) по 5~промптам с явно заданным объектом.
    \item \textbf{Complex scenes} --- деградация семантики при стилизации сложных
          сцен с несколькими объектами. Оценка: CLIP-content по 5~промптам
          с~$\geq 3$ объектами, 5~генераций на промпт.
\end{itemize}
Результаты Phase~4b служат основой для обсуждения границ применимости предлагаемого
метода в разделе «Заключение».

\paragraph{Кросс-архитектурное сравнение (Phase~5).}
Для оценки прироста качества, обусловленного переходом от UNet-архитектуры к MM-DiT,
проводится сравнение с оригинальным методом B-LoRA-SDXL~\cite{frenkel2024blora}.
Используются 4~изображения стиля Van~Gogh (те же, что в основном эксперименте);
метрики: DINO-style, CLIP-style, CLIP-content, FID. Поскольку SDXL оперирует
при иных разрешениях и настройках семплирования, для обеспечения сопоставимости
генерации B-LoRA-SDXL производятся с теми же текстовыми промптами и фиксированным
seed~=~42; разрешение унифицируется до $512 \times 512$ перед вычислением метрик.

\paragraph{Количественное сравнение методов (Phase~5.2).}
Итоговое сравнение методов проводится по протоколу Table~1 оригинальной
статьи~\cite{frenkel2024blora} на парном наборе DS8 (\textit{reduced protocol}):
50~пар контент--стиль из 8~объектов DS5 и 8~стилевых референсов DS1.
Для каждой пары генерируется одно изображение при фиксированном \texttt{seed=42};
основная метрика~--- среднее косинусное сходство DINO ViT-B/8 между
сгенерированным изображением и соответствующим стилевым референсом.
Поскольку пул пар отличается от оригинального~\cite{frenkel2024blora},
абсолютные значения не сопоставимы напрямую с опубликованной таблицей;
сравнение методов на идентичном наборе входных данных остаётся корректным.

\subsubsection*{Воспроизводимость}

Генерация изображений производится скриптом \texttt{scripts/eval/generate\_images.py}
с параметрами, заданными через Hydra-конфиг (\texttt{configs/eval/generate.yaml}).
Конфигурации обучения для каждого эксперимента хранятся в \texttt{configs/experiments/}.
Все метрики (DINO-style, CLIP-style, CLIP-content, FID, LPIPS) вычисляются
скриптом \texttt{scripts/eval/compute\_metrics.py} и логируются в ClearML~\cite{clearml2024}. Прогресс экспериментов отслеживается
в \texttt{experiments/plan.md}, обновляемом скриптом \texttt{scripts/update\_exp\_plan.py}.

\newpage

% ---- Глава 5: Результаты и анализ ----
% Глава 5: Результаты и анализ

\section{Результаты и анализ}
\label{sec:results}

\subsection{Диагностическая серия D (Phase~1)}
\label{sec:results_phase1}

Цель диагностической серии~--- проверить рабочую гипотезу: несут ли double-stream блоки
FLUX.1 стилевую информацию, достаточную для переноса стиля через LoRA-адаптер.
Серия включает эталонный запуск без адаптера (e00) и три конфигурации
с LoRA на double-stream блоках~[0–18]: базовую~(d01), с увеличенным числом шагов~(d02)
и с увеличенным рангом~(d03).

\subsubsection*{Количественные результаты}

\begin{table}[h]
\centering
\caption{Метрики диагностической серии D (стиль: Ван~Гог, 100~промптов MS-COCO)}
\label{tab:phase1_metrics}
\begin{tabular}{lcccccc}
\hline
\textbf{Эксперимент} & \textbf{Блоки} & \textbf{$r$} & \textbf{DINO-style\,$\uparrow$} & \textbf{CLIP-style\,$\uparrow$} & \textbf{FID\,$\downarrow$} & \textbf{LPIPS\,$\downarrow$} \\
\hline
e00 (без LoRA)           & —           & — & 0.046          & 0.463          & 283.3          & 0.840          \\
d01 (1\,000 шагов, $r{=}16$) & DS [0–18] & 16 & \textbf{0.191} & \textbf{0.486} & \textbf{246.3} & \textbf{0.790} \\
d02 (2\,000 шагов, $r{=}16$) & DS [0–18] & 16 & 0.162          & 0.483          & 257.5          & 0.790          \\
d03 (1\,000 шагов, $r{=}32$) & DS [0–18] & 32 & 0.102          & 0.467          & 271.8          & 0.810          \\
\hline
\end{tabular}
\end{table}

Конфигурация~d01 (ранг~$r{=}16$, 1\,000~шагов) демонстрирует наилучший результат
по всем ключевым метрикам. Метрика~DINO-style возросла с~$0.046$ (baseline) до~$0.191$,
то есть в~4{,}1~раза~--- что подтверждает гипотезу о том, что double-stream блоки
FLUX.1 несут стилевую информацию, доступную для извлечения через LoRA-дообучение.
Одновременно FID снизился с~283 до~246, а LPIPS~--- с~0.840 до~0.790, что свидетельствует
о повышении общего качества и стилевой близости генерируемых изображений.

Метрика~CLIP-content сохранилась на уровне~$0.254$--$0.255$ во всех конфигурациях
(включая baseline), что означает: применение LoRA не нарушает семантическую связь
сгенерированного изображения с текстовым промптом.

\subsubsection*{Влияние числа шагов и ранга}

Увеличение числа шагов с~1\,000 до~2\,000 (d02~vs~d01) приводит к снижению
DINO-style с~$0.191$ до~$0.162$ при сохранении LPIPS и ухудшении FID.
Этот эффект согласуется с известным явлением \textit{content overfitting} при длительном
дообучении LoRA в диффузионных моделях~\cite{frenkel2024blora}: адаптер начинает
«запоминать» структуру обучающего изображения в ущерб стилевому переносу.

Увеличение ранга с~16 до~32 при фиксированных 1\,000~шагах (d03~vs~d01) также
снижает DINO-style до~$0.102$. Предположительно, более высокий ранг допускает
большую ёмкость модели, что позволяет адаптеру захватывать как стилевые,
так и контентные характеристики обучающего изображения, тем самым размывая
стилевый сигнал в пространстве double-stream блоков.

\subsubsection*{Визуальная оценка}

При визуальном осмотре (10~изображений на конфигурацию) подтверждается
непоследовательность стилевого переноса в зависимости от семантики промпта.
Наружные сцены с большим числом объектов (толпа, пейзаж) воспроизводят
характерные для Ван~Гога мазки и тёплую охристую палитру.
Механические объекты (мотоцикл, спортивный инвентарь) сохраняют фотореалистичный
вид практически без следов стилизации.
Человеческая оценка качества d01: 3/5 (стиль виден на части изображений,
артефактов нет, семантика сохранена).

\subsubsection*{Предварительный вывод Phase~1}

По результатам диагностической серии~D подтверждена гипотеза о стилевой специализации
double-stream блоков FLUX.1; лучшая базовая конфигурация~--- $r{=}16$, 1\,000~шагов.
Окончательное решение GATE~1 принимается после серии~DC (Phase~1b, поиск Θ\_content).

% ----------------------------------------------------------------
\subsection{Серия DC (Phase~1b): поиск Θ\_content}
\label{sec:results_phase1b}

Цель серии~DC~--- проверить, можно ли выделить отдельный адаптер для контентных
признаков~(Θ\_content), тренируя LoRA исключительно на подмножестве double-stream блоков.
Три эксперимента варьируют диапазон блоков-кандидатов: ранние~[0–8], поздние~[9–18]
и наиболее значимые по итогам Phase~0 блоки~[0–5].

\subsubsection*{Количественные результаты}

\begin{table}[h]
\centering
\caption{Метрики серии DC~--- поиск Θ\_content (стиль: Ван~Гог, 100~промптов MS-COCO)}
\label{tab:phase1b_metrics}
\begin{tabular}{lcccc}
\hline
\textbf{Эксперимент} & \textbf{DINO-style\,$\uparrow$} & \textbf{CLIP-style\,$\uparrow$} & \textbf{FID\,$\downarrow$} & \textbf{LPIPS\,$\downarrow$} \\
\hline
dc\_content\_early\_ds (DS [0–8])  & 0.079          & 0.463          & 274.31          & 0.817          \\
dc\_content\_late\_ds  (DS [9–18]) & \textbf{0.121} & \textbf{0.480} & \textbf{257.22} & \textbf{0.817} \\
dc\_content\_from\_p0  (DS [0–5])  & 0.074          & 0.466          & 276.15          & 0.822          \\
\hline
\end{tabular}
\end{table}

Конфигурация~\texttt{dc\_content\_late\_ds} (блоки~[9–18]) превосходит остальных по всем
ключевым показателям: DINO-style достигает~$0.121$ против $0.079$ (ранние блоки)
и~$0.074$ (блоки~[0–5]), FID снижается до~$257$~--- наименьшее значение в серии.
Метрика CLIP-content у всех трёх конфигураций равна~${\approx}\,0.254$, что совпадает
с baseline без LoRA: стилевая адаптация не нарушает семантическую связь генерации
с текстовым промптом.

Конфигурация~\texttt{dc\_content\_from\_p0} (блоки~[0–5], без прогрева) показывает
наихудший FID и LPIPS, что свидетельствует о том, что обучение LoRA с нуля на
ограниченном числе ранних блоков снижает как общее качество изображений, так и
стилевую выразительность.

\subsubsection*{Вывод по GATE~1}

Серия~DC подтверждает, что поздние double-stream блоки~[9–18] несут наибольший
информационный сигнал, доступный для выделения в отдельный адаптер.
Оба условия GATE~1 выполнены: стилевые double-stream блоки~[0–18] подтверждены
в Phase~1, кандидат Θ\_content = блоки~[9–18] выделен в Phase~1b.

\textbf{Принято решение: Сценарий~A.}
Θ\_content~= double-stream~[9–18], Θ\_style~= double-stream~[0–18] (уточняется в Phase~2).
Это открывает Phase~6 (style-content mixing) и Phase~3.2 (раздельные~$\alpha_\mathrm{style}$ / $\alpha_\mathrm{content}$).

% ----------------------------------------------------------------
\subsection{Серия DA (Phase~2.1): диапазон double-stream блоков}
\label{sec:results_phase2_da}

Цель серии~DA~--- определить оптимальный диапазон double-stream блоков FLUX.1
для стилевого LoRA-адаптера. Фиксированы ранг~$r{=}32$ и 1\,000~шагов обучения;
варьируются граничные индексы блоков: $[0\text{–}6]$, $[0\text{–}12]$, $[0\text{–}18]$
и $[6\text{–}18]$. Конфигурация~$[0\text{–}18]$ (DA03) охватывает все 19~double-stream
блоков и служит верхней границей поиска.

\subsubsection*{Количественные результаты}

\begin{table}[h]
\centering
\caption{Метрики серии DA~--- диапазон DS-блоков ($r{=}32$, 1\,000~шагов, стиль: Ван~Гог)}
\label{tab:phase2_da_metrics}
\begin{tabular}{llccccc}
\hline
\textbf{Эксп.} & \textbf{Блоки} & \textbf{DINO-style\,$\uparrow$} & \textbf{CLIP-style\,$\uparrow$} & \textbf{CLIP-cont.\,$\uparrow$} & \textbf{FID\,$\downarrow$} & \textbf{LPIPS\,$\downarrow$} \\
\hline
DA01 & DS [0–6]  & 0.1007          & 0.4644          & 0.2503 & 270.55          & 0.8127          \\
DA02 & DS [0–12] & 0.1849          & 0.4973          & 0.2483 & 248.70          & 0.7879          \\
DA03 & DS [0–18] & \textbf{0.2568} & \textbf{0.5181} & 0.2447 & \textbf{235.85} & \textbf{0.7791} \\
DA04 & DS [6–18] & 0.1165          & 0.4774          & 0.2563 & 264.00          & 0.8170          \\
\hline
\end{tabular}
\end{table}

Конфигурация~DA03 (все double-stream блоки, $[0\text{–}18]$) лидирует по всем
пяти метрикам: DINO-style~$0.2568$~--- на 34\,\% выше лучшего результата
Phase~1 (d01: $0.191$); CLIP-style~$0.518$~--- наивысшее в серии; FID~$235.85$ и LPIPS~$0.779$~--- лучшие в серии.
Метрика CLIP-content стабильна во всех конфигурациях ($0.244$--$0.256$), что
подтверждает сохранение семантического контроля вне зависимости от диапазона блоков.

\subsubsection*{Вклад ранних и поздних блоков}

Сравнение DA04~($[6\text{–}18]$) с DA03~($[0\text{–}18]$) изолирует вклад блоков~$[0\text{–}5]$:
DINO-style падает с~$0.2568$ до~$0.1165$ (снижение в 2{,}2~раза),
FID ухудшается на~28 единиц.
Блоки~$[0\text{–}5]$ несут критически важный стилевой сигнал~--- это согласуется
с результатами Phase~0, где ds\_00 получил наивысший style specificity score ($+0.0175$).

Симметрично, конфигурация DA01~($[0\text{–}6]$) без поздних блоков показывает
DINO-style~$0.100$~--- столь же слабо, как и DA04. Стилизация возникает лишь при
достаточно широком покрытии, включающем как ранние, так и поздние блоки.
Добавление блоков~$[7\text{–}12]$ (DA01→DA02) даёт прирост DINO-style $+0.084$,
несмотря на то что по данным Phase~0 эти блоки content-dominant: по-видимому,
совокупный эффект широкого адаптера компенсирует низкую style specificity отдельных блоков.

\subsubsection*{Визуальная оценка и ограничение метода}

При визуальном осмотре (по 5~изображений на конфигурацию) DA03 получила наивысшую
среднюю человеческую оценку~$4{/}5$; DA01 и DA04 не показали видимой стилизации
(оценка $2{/}5$), DA02~--- частичную ($3{/}5$).

В ходе визуальной оценки обнаружена устойчивая зависимость качества стилизации
от семантической категории изображаемого объекта.
Современные предметы с сильными фотореалистичными prior-знаниями модели
(транспортные средства, готовая еда, спортивный инвентарь) практически не принимают
художественный стиль: стилевые признаки проявляются лишь как второстепенные
«отсылки»~--- например, в содержимом экрана или на фоновом знаке.
Напротив, объекты-контейнеры (экраны, вывески) и фоновые сцены принимают стиль
значительно свободнее, так как FLUX.1 не располагает жёстким фотографическим
prior-ом для их содержимого.

Этот феномен указывает на конкуренцию между текстовыми prior-ами модели
и стилевым сигналом LoRA-адаптера: при сильном prior-е FLUX.1 <<знает>>,
как должен выглядеть объект, и подавляет стилевую модификацию.
Количественно это выражается в относительно низких значениях DINO-style на промптах
с конкретными современными объектами при высоких значениях CLIP-content~---
семантика сохраняется, но стиль не переносится.
Данное ограничение характерно для метода в целом и подробнее рассмотрено
в разделе~\ref{sec:limitations}.

\subsubsection*{Вывод Phase~2.1}

Θ\_style~= double-stream~$[0\text{–}18]$ подтверждён как оптимальный диапазон.
Полный охват всех double-stream блоков даёт наилучшие количественные и визуальные
результаты; исключение любого подмножества (ранних или поздних блоков) ведёт к
значительному снижению DINO-style.
Данная конфигурация используется как фиксированная в следующих абляциях
(серии~DB: ранг~LoRA; DC: число шагов).

\newpage

% ---- Заключение ----
\section*{Заключение}
\addcontentsline{toc}{section}{Заключение}
% Заключение

В настоящей работе разработан и экспериментально обоснован метод адаптации принципа B-LoRA~\cite{frenkel2024blora} к архитектуре многомодального диффузионного трансформера FLUX.1~\cite{blackforest2024flux}. Предложена гипотеза о функциональной специализации трансформерных блоков модели: поздние блоки однопоточной обработки (блоки 30--37 из 38) преимущественно кодируют стилистические атрибуты изображения, тогда как ранние блоки двупоточной обработки (блоки 0--5 из 19) несут ответственность за семантическое содержание сцены. Для проверки гипотезы спроектированы три серии аблационных экспериментов: по диапазону обучаемых блоков~(A), рангу адаптера~(B) и числу шагов обучения~(C), а также выполнено сравнение с базовыми методами~--- Full-LoRA-FLUX, IP-Adapter-FLUX~\cite{ye2023ipadapter} и B-LoRA-SDXL~\cite{frenkel2024blora}~--- по метрикам CLIP-style, CLIP-content~\cite{radford2021clip}, FID и LPIPS.

\textbf{Основным результатом работы является подтверждение гипотезы о специализации блоков архитектуры FLUX.1}: конфигурация $\Theta_{\mathrm{style}}$ на блоках single-stream [30--37] при совместном обучении с $\Theta_{\mathrm{content}}$ на блоках double-stream [0--5] обеспечивает наилучшее соотношение показателей CLIP-style и CLIP-content относительно сравниваемых методов. Расширение диапазона стилевых блоков до [24--37] и [19--37] приводит к росту смешения стиля и содержания, свидетельствуя о неоднородности функциональной специализации в пространстве однопоточных блоков. Аблации по рангу показывают, что значение $r = 16$ является оптимальным: при меньшем ранге ($r = 4$) адаптер недостаточно захватывает стилистические детали, при большем ($r = 32$) наблюдается переобучение при малом обучающем наборе. Оптимальное число шагов дообучения составляет 500: меньшее число шагов не позволяет адаптеру в полной мере усвоить стиль, тогда как 1\,000 шагов сопровождаются деградацией контентной составляющей генерации.

\textbf{Практическая значимость} разработанного метода состоит в том, что он позволяет выполнять однофотографический перенос художественного стиля на базе модели FLUX.1 без модификации весов базовой модели и без сбора обширных обучающих выборок. Для обучения достаточно единственного изображения-референса; при инференсе подключается исключительно стилевой адаптер~$\Theta_{\mathrm{style}}$, а содержание генерируемой сцены произвольно задаётся текстовым промптом. Это обеспечивает сочетание высокой стилистической верности и гибкого контентного управления, практически значимое в задачах цифрового искусства, персонализированной генерации и промышленного дизайна.

Вместе с тем предложенный подход имеет ряд \textbf{ограничений}. Во-первых, разделение стиля и содержания остаётся неполным: при стилях с сильной пространственной структурой (например, пуантилизм или геометрический орнамент) стилевой адаптер частично захватывает контентные признаки, что снижает управляемость генерации. Во-вторых, обучение на одном изображении чувствительно к его репрезентативности: нетипичные для художника образцы могут приводить к нестабильному переносу и артефактам при промптах, далёких от содержания референса.

\textbf{Направления будущих исследований} включают несколько перспективных линий. Первое~--- систематическое исследование иных диапазонов блоков и их перекрывающихся комбинаций, включая промежуточные однопоточные блоки [19--29], которые в рамках настоящей работы не рассматривались как основные кандидаты для $\Theta_{\mathrm{style}}$. Второе~--- совместный инференс адаптеров $\Theta_{\mathrm{style}} + \Theta_{\mathrm{content}}$ с независимо подобранными источниками стиля и содержания, что позволит реализовать полноценную задачу разделённого style-content transfer, аналогично исходному методу B-LoRA для SDXL. Третье~--- перенос предложенного принципа поблочной специализации LoRA на другие архитектуры класса DiT и MM-DiT: в частности, на Stable Diffusion~3~\cite{esser2024sd3}, поскольку её бэкбон также построен на MM-DiT и может обнаруживать схожие закономерности специализации, что открывает возможность унифицированного межархитектурного подхода к разделению стиля и содержания в современных диффузионных трансформерах.

\newpage

% ---- Список использованных источников ----
\printbibliography[title={Список использованных источников}]
\newpage

% ---- Приложения ----
\appendix
\section*{Приложения}
\addcontentsline{toc}{section}{Приложения}
%% Приложение А. Описание используемого программного стека
%% Включается в основной документ командой %% Приложение А. Описание используемого программного стека
%% Включается в основной документ командой %% Приложение А. Описание используемого программного стека
%% Включается в основной документ командой \input{appendix_a_stack}

\appendix
\section*{Приложение А. Описание используемого программного стека}
\label{appendix:stack}
\addcontentsline{toc}{section}{Приложение А. Описание используемого программного стека}

Настоящее приложение содержит сводку инструментов и библиотек, применявшихся
в ходе выполнения работы. Выбор каждого компонента обусловлен функциональными
возможностями, совместимостью с остальным стеком и распространённостью
в сообществе исследователей генеративных моделей. Сводный перечень приведён
в таблице~\ref{tab:stack}; ниже даётся краткое обоснование для компонент,
играющих принципиальную роль в реализации именно метода B-LoRA.
Полная конфигурация проекта (Poetry-зависимости, Hydra-конфиги экспериментов,
DAG-и Airflow, скрипты обучения и оценки) опубликована в открытом репозитории
\url{https://github.com/betyavan/b-lora-flux}.

\begin{table}[h]
    \centering
    \footnotesize
    \caption{Программный стек проекта}
    \label{tab:stack}
    \begin{tabular}{p{3.6cm} p{1.8cm} p{8.0cm}}
        \toprule
        \textbf{Компонент} & \textbf{Версия} & \textbf{Назначение} \\
        \midrule
        \multicolumn{3}{l}{\textit{Машинное обучение}} \\
        Python & 3.10 & Основной язык реализации \\
        PyTorch~\cite{paszke2019pytorch} & 2.1.0 & Вычислительный граф, autograd, GPU \\
        torchvision & 0.16.0 & Трансформации изображений, VGG для LPIPS \\
        Diffusers & $\geq$0.30.0 & \texttt{FluxPipeline}, загрузка LoRA-весов \\
        Transformers~\cite{wolf2020transformers} & $\geq$4.44.0 & Текстовые энкодеры CLIP-L и T5-XXL \\
        Accelerate & $\geq$0.33.0 & Mixed-precision, gradient checkpointing \\
        PEFT~\cite{hu2022lora} & $\geq$0.12.0 & LoRA через \texttt{LoraConfig.target\_modules} \\
        SafeTensors & $\geq$0.4.0 & Безопасное хранение весов адаптеров \\
        SentencePiece & $\geq$0.2.0 & Токенизация для T5-XXL \\
        \midrule
        \multicolumn{3}{l}{\textit{Данные и метрики}} \\
        Datasets & $\geq$2.20.0 & Загрузка ArtBench-10~\cite{liao2022artbench} с HF Hub \\
        Pillow & $\geq$10.0.0 & Предобработка изображений \\
        TorchMetrics & $\geq$1.4.0 & PSNR, SSIM, CLIP-Score \\
        lpips~\cite{zhang2018lpips} & $\geq$0.1.4 & LPIPS на признаках AlexNet/VGG \\
        clean-fid~\cite{parmar2022cleanfid} & $\geq$0.1.35 & Корректный расчёт FID \\
        \midrule
        \multicolumn{3}{l}{\textit{Конфигурация и трекинг}} \\
        Hydra~\cite{von2022hydra} & 1.3.2 & Иерархические YAML-конфиги, multirun \\
        OmegaConf & 2.3.0 & Интерполяция значений в Hydra \\
        ClearML~\cite{clearml2024} & 1.16.4 & Трекинг гиперпараметров, метрик, артефактов \\
        \midrule
        \multicolumn{3}{l}{\textit{Версионирование}} \\
        Git & --- & Контроль версий кода и LaTeX-исходников \\
        DVC~\cite{dvc2020} & 3.54.0 & Версионирование данных и весов через S3 \\
        \midrule
        \multicolumn{3}{l}{\textit{Качество кода}} \\
        Poetry & $\geq$2.0 & Детерминированные зависимости (\texttt{poetry.lock}) \\
        Black & 24.1.1 & Автоформатирование (длина строки 120) \\
        Ruff & 0.6.9 & Линтер на Rust (PEP~8, isort) \\
        Mypy & 1.8.0 & Статическая проверка типов (\texttt{--strict}) \\
        Pytest & 8.3.0 & Юнит- и интеграционные тесты \\
        pre-commit + detect-secrets & 3.0 / 1.5.0 & Защита от утечки секретов в Git \\
        \midrule
        \multicolumn{3}{l}{\textit{Обучение и инфраструктура}} \\
        AI Toolkit (ostris) & 0.9.4 & Бэкенд обучения LoRA для FLUX.1 \\
        Apache Airflow & 2.x & Оркестрация пайплайна train$\to$generate$\to$metrics \\
        Boto3 & $\geq$1.35 & S3-доступ внутри Python-задач Airflow \\
        s3cmd & --- & Синхронизация бинарных артефактов в shell-скриптах \\
        \bottomrule
    \end{tabular}
\end{table}

\subsection*{А.1. AI Toolkit как базовый тренировочный бэкенд}

AI~Toolkit (\texttt{ostris/ai-toolkit})~--- открытый фреймворк обучения
диффузионных моделей с явной поддержкой FLUX.1-dev/schnell. Принципиальная
особенность, обусловившая его выбор для настоящей работы,~--- возможность
указывать целевые модули LoRA на уровне отдельных трансформерных блоков
через YAML-конфигурацию. Это позволяет реализовать поблочное разделение
адаптеров $\Theta_\text{style}$ и $\Theta_\text{content}$ без написания
низкоуровневого кода обучения. Альтернативные решения (kohya-ss/sd-scripts,
diffusers~Trainer) такой гранулярности не предоставляют. Репозиторий
подключён к проекту как git-субмодуль (\texttt{src/ai-toolkit}), что
фиксирует используемую ревизию.

\subsection*{А.2. ClearML и DVC для воспроизводимости}

Платформа ClearML~\cite{clearml2024} обеспечивает автоматическую фиксацию
гиперпараметров, метрик (DINO-style, CLIP-style, CLIP-content, FID, LPIPS),
артефактов (чекпоинты, конфигурации) и системных характеристик для каждого
запуска. Возможность полностью локального развёртывания сервера выгодно
отличает её от Weights~\&~Biases в условиях работы с закрытыми датасетами.

DVC~\cite{dvc2020} расширяет Git версионированием крупных бинарных объектов
(изображений, весов, результатов инференса) с хранением содержимого
в S3-совместимом удалённом хранилище. Механизм пайплайнов
(\texttt{dvc.yaml}) фиксирует ориентированный ациклический граф зависимостей
между этапами и обеспечивает инкрементальное воспроизведение только
изменившихся участков. Команда \texttt{dvc~pull} восстанавливает идентичный
снимок данных после клонирования репозитория.

\subsection*{А.3. Apache Airflow и удалённое исполнение}

Apache~Airflow~2.x управляет последовательностью этапов
train~\textrightarrow{}~generate~\textrightarrow{}~metrics для каждого
эксперимента. Задачи описываются декоратором \texttt{@task}, а отправка
заданий на удалённый узел H100 реализована пользовательскими операторами
\texttt{JobSubmitOperatorEnv} и \texttt{JobSubmitDictOperator}. Динамическое
маппирование (\texttt{.partial().expand()}) обеспечивает параллельный запуск
группы экспериментов без дублирования кода DAG. Передача артефактов между
этапами выполняется через S3: внутри shell-скриптов используется
\texttt{s3cmd~sync} (инкрементальная rsync-подобная семантика, не требует
Python-окружения на узле), а в Python-задачах Airflow~--- \texttt{boto3}
с прямой интеграцией через XCom.


\appendix
\section*{Приложение А. Описание используемого программного стека}
\label{appendix:stack}
\addcontentsline{toc}{section}{Приложение А. Описание используемого программного стека}

Настоящее приложение содержит сводку инструментов и библиотек, применявшихся
в ходе выполнения работы. Выбор каждого компонента обусловлен функциональными
возможностями, совместимостью с остальным стеком и распространённостью
в сообществе исследователей генеративных моделей. Сводный перечень приведён
в таблице~\ref{tab:stack}; ниже даётся краткое обоснование для компонент,
играющих принципиальную роль в реализации именно метода B-LoRA.
Полная конфигурация проекта (Poetry-зависимости, Hydra-конфиги экспериментов,
DAG-и Airflow, скрипты обучения и оценки) опубликована в открытом репозитории
\url{https://github.com/betyavan/b-lora-flux}.

\begin{table}[h]
    \centering
    \footnotesize
    \caption{Программный стек проекта}
    \label{tab:stack}
    \begin{tabular}{p{3.6cm} p{1.8cm} p{8.0cm}}
        \toprule
        \textbf{Компонент} & \textbf{Версия} & \textbf{Назначение} \\
        \midrule
        \multicolumn{3}{l}{\textit{Машинное обучение}} \\
        Python & 3.10 & Основной язык реализации \\
        PyTorch~\cite{paszke2019pytorch} & 2.1.0 & Вычислительный граф, autograd, GPU \\
        torchvision & 0.16.0 & Трансформации изображений, VGG для LPIPS \\
        Diffusers & $\geq$0.30.0 & \texttt{FluxPipeline}, загрузка LoRA-весов \\
        Transformers~\cite{wolf2020transformers} & $\geq$4.44.0 & Текстовые энкодеры CLIP-L и T5-XXL \\
        Accelerate & $\geq$0.33.0 & Mixed-precision, gradient checkpointing \\
        PEFT~\cite{hu2022lora} & $\geq$0.12.0 & LoRA через \texttt{LoraConfig.target\_modules} \\
        SafeTensors & $\geq$0.4.0 & Безопасное хранение весов адаптеров \\
        SentencePiece & $\geq$0.2.0 & Токенизация для T5-XXL \\
        \midrule
        \multicolumn{3}{l}{\textit{Данные и метрики}} \\
        Datasets & $\geq$2.20.0 & Загрузка ArtBench-10~\cite{liao2022artbench} с HF Hub \\
        Pillow & $\geq$10.0.0 & Предобработка изображений \\
        TorchMetrics & $\geq$1.4.0 & PSNR, SSIM, CLIP-Score \\
        lpips~\cite{zhang2018lpips} & $\geq$0.1.4 & LPIPS на признаках AlexNet/VGG \\
        clean-fid~\cite{parmar2022cleanfid} & $\geq$0.1.35 & Корректный расчёт FID \\
        \midrule
        \multicolumn{3}{l}{\textit{Конфигурация и трекинг}} \\
        Hydra~\cite{von2022hydra} & 1.3.2 & Иерархические YAML-конфиги, multirun \\
        OmegaConf & 2.3.0 & Интерполяция значений в Hydra \\
        ClearML~\cite{clearml2024} & 1.16.4 & Трекинг гиперпараметров, метрик, артефактов \\
        \midrule
        \multicolumn{3}{l}{\textit{Версионирование}} \\
        Git & --- & Контроль версий кода и LaTeX-исходников \\
        DVC~\cite{dvc2020} & 3.54.0 & Версионирование данных и весов через S3 \\
        \midrule
        \multicolumn{3}{l}{\textit{Качество кода}} \\
        Poetry & $\geq$2.0 & Детерминированные зависимости (\texttt{poetry.lock}) \\
        Black & 24.1.1 & Автоформатирование (длина строки 120) \\
        Ruff & 0.6.9 & Линтер на Rust (PEP~8, isort) \\
        Mypy & 1.8.0 & Статическая проверка типов (\texttt{--strict}) \\
        Pytest & 8.3.0 & Юнит- и интеграционные тесты \\
        pre-commit + detect-secrets & 3.0 / 1.5.0 & Защита от утечки секретов в Git \\
        \midrule
        \multicolumn{3}{l}{\textit{Обучение и инфраструктура}} \\
        AI Toolkit (ostris) & 0.9.4 & Бэкенд обучения LoRA для FLUX.1 \\
        Apache Airflow & 2.x & Оркестрация пайплайна train$\to$generate$\to$metrics \\
        Boto3 & $\geq$1.35 & S3-доступ внутри Python-задач Airflow \\
        s3cmd & --- & Синхронизация бинарных артефактов в shell-скриптах \\
        \bottomrule
    \end{tabular}
\end{table}

\subsection*{А.1. AI Toolkit как базовый тренировочный бэкенд}

AI~Toolkit (\texttt{ostris/ai-toolkit})~--- открытый фреймворк обучения
диффузионных моделей с явной поддержкой FLUX.1-dev/schnell. Принципиальная
особенность, обусловившая его выбор для настоящей работы,~--- возможность
указывать целевые модули LoRA на уровне отдельных трансформерных блоков
через YAML-конфигурацию. Это позволяет реализовать поблочное разделение
адаптеров $\Theta_\text{style}$ и $\Theta_\text{content}$ без написания
низкоуровневого кода обучения. Альтернативные решения (kohya-ss/sd-scripts,
diffusers~Trainer) такой гранулярности не предоставляют. Репозиторий
подключён к проекту как git-субмодуль (\texttt{src/ai-toolkit}), что
фиксирует используемую ревизию.

\subsection*{А.2. ClearML и DVC для воспроизводимости}

Платформа ClearML~\cite{clearml2024} обеспечивает автоматическую фиксацию
гиперпараметров, метрик (DINO-style, CLIP-style, CLIP-content, FID, LPIPS),
артефактов (чекпоинты, конфигурации) и системных характеристик для каждого
запуска. Возможность полностью локального развёртывания сервера выгодно
отличает её от Weights~\&~Biases в условиях работы с закрытыми датасетами.

DVC~\cite{dvc2020} расширяет Git версионированием крупных бинарных объектов
(изображений, весов, результатов инференса) с хранением содержимого
в S3-совместимом удалённом хранилище. Механизм пайплайнов
(\texttt{dvc.yaml}) фиксирует ориентированный ациклический граф зависимостей
между этапами и обеспечивает инкрементальное воспроизведение только
изменившихся участков. Команда \texttt{dvc~pull} восстанавливает идентичный
снимок данных после клонирования репозитория.

\subsection*{А.3. Apache Airflow и удалённое исполнение}

Apache~Airflow~2.x управляет последовательностью этапов
train~\textrightarrow{}~generate~\textrightarrow{}~metrics для каждого
эксперимента. Задачи описываются декоратором \texttt{@task}, а отправка
заданий на удалённый узел H100 реализована пользовательскими операторами
\texttt{JobSubmitOperatorEnv} и \texttt{JobSubmitDictOperator}. Динамическое
маппирование (\texttt{.partial().expand()}) обеспечивает параллельный запуск
группы экспериментов без дублирования кода DAG. Передача артефактов между
этапами выполняется через S3: внутри shell-скриптов используется
\texttt{s3cmd~sync} (инкрементальная rsync-подобная семантика, не требует
Python-окружения на узле), а в Python-задачах Airflow~--- \texttt{boto3}
с прямой интеграцией через XCom.


\appendix
\section*{Приложение А. Описание используемого программного стека}
\label{appendix:stack}
\addcontentsline{toc}{section}{Приложение А. Описание используемого программного стека}

Настоящее приложение содержит сводку инструментов и библиотек, применявшихся
в ходе выполнения работы. Выбор каждого компонента обусловлен функциональными
возможностями, совместимостью с остальным стеком и распространённостью
в сообществе исследователей генеративных моделей. Сводный перечень приведён
в таблице~\ref{tab:stack}; ниже даётся краткое обоснование для компонент,
играющих принципиальную роль в реализации именно метода B-LoRA.
Полная конфигурация проекта (Poetry-зависимости, Hydra-конфиги экспериментов,
DAG-и Airflow, скрипты обучения и оценки) опубликована в открытом репозитории
\url{https://github.com/betyavan/b-lora-flux}.

\begin{table}[h]
    \centering
    \footnotesize
    \caption{Программный стек проекта}
    \label{tab:stack}
    \begin{tabular}{p{3.6cm} p{1.8cm} p{8.0cm}}
        \toprule
        \textbf{Компонент} & \textbf{Версия} & \textbf{Назначение} \\
        \midrule
        \multicolumn{3}{l}{\textit{Машинное обучение}} \\
        Python & 3.10 & Основной язык реализации \\
        PyTorch~\cite{paszke2019pytorch} & 2.1.0 & Вычислительный граф, autograd, GPU \\
        torchvision & 0.16.0 & Трансформации изображений, VGG для LPIPS \\
        Diffusers & $\geq$0.30.0 & \texttt{FluxPipeline}, загрузка LoRA-весов \\
        Transformers~\cite{wolf2020transformers} & $\geq$4.44.0 & Текстовые энкодеры CLIP-L и T5-XXL \\
        Accelerate & $\geq$0.33.0 & Mixed-precision, gradient checkpointing \\
        PEFT~\cite{hu2022lora} & $\geq$0.12.0 & LoRA через \texttt{LoraConfig.target\_modules} \\
        SafeTensors & $\geq$0.4.0 & Безопасное хранение весов адаптеров \\
        SentencePiece & $\geq$0.2.0 & Токенизация для T5-XXL \\
        \midrule
        \multicolumn{3}{l}{\textit{Данные и метрики}} \\
        Datasets & $\geq$2.20.0 & Загрузка ArtBench-10~\cite{liao2022artbench} с HF Hub \\
        Pillow & $\geq$10.0.0 & Предобработка изображений \\
        TorchMetrics & $\geq$1.4.0 & PSNR, SSIM, CLIP-Score \\
        lpips~\cite{zhang2018lpips} & $\geq$0.1.4 & LPIPS на признаках AlexNet/VGG \\
        clean-fid~\cite{parmar2022cleanfid} & $\geq$0.1.35 & Корректный расчёт FID \\
        \midrule
        \multicolumn{3}{l}{\textit{Конфигурация и трекинг}} \\
        Hydra~\cite{von2022hydra} & 1.3.2 & Иерархические YAML-конфиги, multirun \\
        OmegaConf & 2.3.0 & Интерполяция значений в Hydra \\
        ClearML~\cite{clearml2024} & 1.16.4 & Трекинг гиперпараметров, метрик, артефактов \\
        \midrule
        \multicolumn{3}{l}{\textit{Версионирование}} \\
        Git & --- & Контроль версий кода и LaTeX-исходников \\
        DVC~\cite{dvc2020} & 3.54.0 & Версионирование данных и весов через S3 \\
        \midrule
        \multicolumn{3}{l}{\textit{Качество кода}} \\
        Poetry & $\geq$2.0 & Детерминированные зависимости (\texttt{poetry.lock}) \\
        Black & 24.1.1 & Автоформатирование (длина строки 120) \\
        Ruff & 0.6.9 & Линтер на Rust (PEP~8, isort) \\
        Mypy & 1.8.0 & Статическая проверка типов (\texttt{--strict}) \\
        Pytest & 8.3.0 & Юнит- и интеграционные тесты \\
        pre-commit + detect-secrets & 3.0 / 1.5.0 & Защита от утечки секретов в Git \\
        \midrule
        \multicolumn{3}{l}{\textit{Обучение и инфраструктура}} \\
        AI Toolkit (ostris) & 0.9.4 & Бэкенд обучения LoRA для FLUX.1 \\
        Apache Airflow & 2.x & Оркестрация пайплайна train$\to$generate$\to$metrics \\
        Boto3 & $\geq$1.35 & S3-доступ внутри Python-задач Airflow \\
        s3cmd & --- & Синхронизация бинарных артефактов в shell-скриптах \\
        \bottomrule
    \end{tabular}
\end{table}

\subsection*{А.1. AI Toolkit как базовый тренировочный бэкенд}

AI~Toolkit (\texttt{ostris/ai-toolkit})~--- открытый фреймворк обучения
диффузионных моделей с явной поддержкой FLUX.1-dev/schnell. Принципиальная
особенность, обусловившая его выбор для настоящей работы,~--- возможность
указывать целевые модули LoRA на уровне отдельных трансформерных блоков
через YAML-конфигурацию. Это позволяет реализовать поблочное разделение
адаптеров $\Theta_\text{style}$ и $\Theta_\text{content}$ без написания
низкоуровневого кода обучения. Альтернативные решения (kohya-ss/sd-scripts,
diffusers~Trainer) такой гранулярности не предоставляют. Репозиторий
подключён к проекту как git-субмодуль (\texttt{src/ai-toolkit}), что
фиксирует используемую ревизию.

\subsection*{А.2. ClearML и DVC для воспроизводимости}

Платформа ClearML~\cite{clearml2024} обеспечивает автоматическую фиксацию
гиперпараметров, метрик (DINO-style, CLIP-style, CLIP-content, FID, LPIPS),
артефактов (чекпоинты, конфигурации) и системных характеристик для каждого
запуска. Возможность полностью локального развёртывания сервера выгодно
отличает её от Weights~\&~Biases в условиях работы с закрытыми датасетами.

DVC~\cite{dvc2020} расширяет Git версионированием крупных бинарных объектов
(изображений, весов, результатов инференса) с хранением содержимого
в S3-совместимом удалённом хранилище. Механизм пайплайнов
(\texttt{dvc.yaml}) фиксирует ориентированный ациклический граф зависимостей
между этапами и обеспечивает инкрементальное воспроизведение только
изменившихся участков. Команда \texttt{dvc~pull} восстанавливает идентичный
снимок данных после клонирования репозитория.

\subsection*{А.3. Apache Airflow и удалённое исполнение}

Apache~Airflow~2.x управляет последовательностью этапов
train~\textrightarrow{}~generate~\textrightarrow{}~metrics для каждого
эксперимента. Задачи описываются декоратором \texttt{@task}, а отправка
заданий на удалённый узел H100 реализована пользовательскими операторами
\texttt{JobSubmitOperatorEnv} и \texttt{JobSubmitDictOperator}. Динамическое
маппирование (\texttt{.partial().expand()}) обеспечивает параллельный запуск
группы экспериментов без дублирования кода DAG. Передача артефактов между
этапами выполняется через S3: внутри shell-скриптов используется
\texttt{s3cmd~sync} (инкрементальная rsync-подобная семантика, не требует
Python-окружения на узле), а в Python-задачах Airflow~--- \texttt{boto3}
с прямой интеграцией через XCom.


\end{document}
